\documentclass[acmsmall,screen,nonacm]{acmart}

\usepackage{bcprules}
\usepackage{mathtools}
\usepackage[nameinlink,noabbrev]{cleveref}

\settopmatter{printacmref=false,printccs=false,printfolios=true}
\raggedbottom

\newtheorem{theorem}{Theorem}
\newtheorem{lemma}[theorem]{Lemma}

\newcounter{RuleRef}
\newcommand{\ruledef}[1]{%
  \refstepcounter{RuleRef}\textsc{#1}\label{rule:#1}}
\newcommand{\ruleuse}[1]{%
  \hyperref[rule:#1]{\textsc{#1}}}
\newcommand{\ruleref}[1]{%
  \textsc{(\hyperref[rule:#1]{#1})}}

\newcommand{\lambdap}{\ensuremath{\lambda_{\mathsf p}}}
\newcommand{\ctxempty}{\varnothing}
\newcommand{\sing}[1]{\{#1\}}
\newcommand{\pairty}[4]{\langle #1\!:\!#2;\,#3\!:\!#4\rangle}
\newcommand{\pairval}[3]{\langle #1;\,#2=#3\rangle}
\newcommand{\subst}[3]{#1[#2/#3]}
\newcommand{\wf}{\mathrel{\mathsf{wf}}}
\newcommand{\final}{\mathsf{final}}
\newcommand{\steps}{\longrightarrow^{*}}
\newcommand{\sstep}{\longrightarrow}
\newcommand{\resolve}{\Downarrow}
\newcommand{\loc}[1]{\mathsf{loc}(#1)}
\newcommand{\typeref}[1]{\mathsf{type}(#1)}
\newcommand{\referentof}[1]{\mathsf{ref}(#1)}
\newcommand{\realizes}{\Vdash}
\newcommand{\coerce}{\mathrel{\rightsquigarrow}}
\newcommand{\bodycl}[4]{\mathsf{Body}_{#1}(#2,#3,#4)}
\newcommand{\membercl}[5]{%
  \mathsf{Member}_{#1}(#2\!:\!#3;\,#4,#5)}
\newcommand{\judgheader}[3][]{%
  \par\addvspace{0.75em}%
  \noindent\textbf{#2}%
  \if\relax\detokenize{#1}\relax\else\quad #1\fi
  \hfill{\setlength{\fboxsep}{4pt}\setlength{\fboxrule}{0.5pt}%
    \fbox{\(\strut #3\)}}\par
  \nobreak\addvspace{0.85em}%
}
\newcommand{\plateheader}[2]{%
  \par\addvspace{1.1em}%
  \noindent\textbf{#1}%
  \if\relax\detokenize{#2}\relax\else
    \hfill{\setlength{\fboxsep}{3pt}\setlength{\fboxrule}{0.5pt}%
      \fbox{\(\strut #2\)}}%
  \fi\par
  \nobreak\addvspace{0.65em}%
}
\newcommand{\rulesetup}{%
  \small
  \smallrulenames
  \setlength{\afterruleskip}{2em plus 0.2em minus 0.1em}%
  \setlength{\labelminsep}{0.6em}%
  \renewcommand{\arraystretch}{1.15}%
  \renewcommand{\andalso}{\qquad}%
}

\title{Semantic Coercions for Paths through Dependent Pairs}
\author{Oliver Bra\v{c}evac}
\author{Martin Odersky}
\renewcommand{\shortauthors}{Bra\v{c}evac and Odersky}
\authorsaddresses{}

\begin{document}

\begin{abstract}
We study a small call-by-value calculus of path-dependent types in which
paths traverse immutable dependent pairs. Pairs carry term members or
abstract type members, and their second component may mention the first. The
central subtyping rule is covariant in both components: it compares dependent
members under the source first-component type. This rule resists the usual
syntax-directed preservation argument because its member premise must be
instantiated at a concrete store location that is unavailable in the source
derivation.

The soundness proof interprets each subtyping derivation as a finite,
store-indexed semantic coercion. A coercion maps realizations of its source
type to realizations of its target type. Function codomains and pair members
are retained as closures until execution supplies the corresponding
location. For a pair coercion, the saved member derivation is instantiated at
the first component stored in the pair and then used to reinterpret the
stored member. The action is well founded because every location named by a
stored pair precedes the pair cell in the allocation order; a stored type is
a base case. We give the complete
calculus, operational semantics, coercion construction, machine invariant,
and progress and preservation theorems. The development is mechanized in
Lean 4.
\end{abstract}

\keywords{path-dependent types, dependent pairs, subtyping, semantic
coercions, type soundness, Lean}

\maketitle

\section{Introduction}

Path-dependent types associate type information with stable term paths. DOT
calculi place this mechanism in a rich object calculus with recursive self
types, refinements, intersections, and abstract type members
\citep{amin2014foundations,rompf2016dot}. The interaction among those
features has led to several distinct soundness arguments. The calculus in
this paper, \lambdap{}, isolates a smaller question: how much of paths and
abstract type selection can be based on ordinary dependent pairs?

A \lambdap{} value is an abstraction or an immutable pair
\(\pairval{y}{a}{\delta}\). Its first component is a store location. The
second component is either another location or a type definition, and its
signature may refer to the first component. Paths consist of variables,
first projections, and labelled selections. A failed selection continues
through the first component, allowing nested pairs to represent extensible
records. Term singleton types \(\sing{p}\) record path identity. Abstract
selections \(p.A\) range between a stored lower and upper bound. The language
also contains dependent functions, declarative subtyping with primitive
transitivity, and explicit consistency premises for interval bounds.

Write \(\Gamma\vdash\tau<:\tau'\) when signature \(\tau\) is a subtype of
\(\tau'\) under typing context \(\Gamma\). The difficult rule is covariance
of dependent pairs. Here \(S,S'\) range over proper types, while
\(\tau,\tau'\) range over member signatures, which may be either proper
types or type intervals. The type \(\pairty{x}{S}{a}{\tau}\) binds \(x\) in
\(\tau\), where \(x\) denotes the first component and \(a\) is the member
label:
\begingroup
\rulesetup
\typicallabel{s-pair}
\infrule[\ruleuse{s-pair}]
  {\Gamma\vdash S<:S'
   \andalso
   \Gamma,x:S\vdash \tau<:\tau'}
  {\Gamma\vdash
    \pairty{x}{S}{a}{\tau}<:\pairty{x}{S'}{a}{\tau'}}
\endgroup
The source pair supplies a component of type \(S\), so \(S\) is the natural
assumption for comparing the members. At run time, however, the member
signature is opened with the concrete location \(y\) stored in the pair. A
proof of preservation must use the second premise at that location and must
relate source paths to the referents found by run-time resolution.

We make this use of subtyping evidence explicit in the metatheory. Fix a
store \(\sigma\) and an environment mapping source variables to locations in
that store. A path's \emph{referent} is either a location or a stored type. A
subtyping derivation then yields a semantic coercion
\[
  \sigma\vdash \tau\coerce \tau'.
\]
Its meaning is a closure property of semantic typing: whenever a referent
realizes \(\tau\), the same referent realizes \(\tau'\). A coercion is finite
evidence used to prove this property.
Transitivity becomes coercion composition. Rules for singleton and
abstract selection retain the path resolutions and interval evidence needed
by their actions. Function and pair rules retain dependent premises until a
concrete binder location is known.

For \ruleref{s-pair}, construction produces a coercion for \(S<:S'\) and a
closure containing the derivation of \(\tau<:\tau'\). Acting on a pair at location
\(z\) exposes its first component \(y\) and its stored definition \(\delta\).
The first coercion maps the realization of \(y\) from \(S\) to \(S'\). The
closure extends its saved environment by \(x\mapsto y\), producing a
coercion from \(\tau[y/x]\) to \(\tau'[y/x]\), which acts on the referent of
\(\delta\); the notation replaces the formal parameter \(x\) by \(y\).
This dynamically produced coercion is not a structural child of
the pair coercion. The first component \(y\) and any location stored in
\(\delta\) were allocated before \(z\). Coercion action uses a lexicographic
measure consisting of the referent's allocation rank and coercion-tree size.

The proof uses these coercions only as semantic evidence. Programs and
machine states contain the source terms given in \cref{sec:calculus};
execution performs no casts or type checks. Typed intermediate calculi such
as System FC instead place equality evidence and casts in the intermediate
program and erase them before execution \citep{sulzmann2007systemfc}.

The contributions are:
\begin{itemize}
\item a small calculus combining arbitrary stable paths, term singletons,
  abstract type selections, dependent functions, and dependent pairs with
  proper or interval members;
\item a store-indexed coercion interpretation of its full declarative
  subtyping relation, including primitive transitivity and unrestricted
  dependent-pair covariance;
\item a well-founded action for pair coercions based on the allocation order
  of immutable stores; and
\item a Lean 4 mechanization of path resolution, coercion construction and
  action, a syntax-directed machine invariant, progress, preservation, and
  non-stuckness of every finite execution prefix.
\end{itemize}

The result is type safety for the calculus as presented. It does not establish
normalization, decidability of subtyping, or admissibility of transitivity.

\section{The calculus}
\label{sec:calculus}

\lambdap{} isolates path-dependent types in a call-by-value calculus with
immutable dependent pairs. Terms are in monadic normal form: applications
take paths as operands, and a pair contains a location as its first component
together with either another location or a type definition. The operational
semantics combines path resolution with a CK machine.

\subsection{Source language}
\label{sec:syntax}

Proper types and member signatures are defined mutually. In the grammar,
\(S,T,U\) range over proper types; \(\tau,\tau'\) range over proper
or interval signatures; \(p,q\) range over paths; \(t,u\) range over
terms; \(v\) ranges over values; \(x,y,z\) range over variables; \(a,b\)
range over labels; \(\delta\) ranges over member definitions; and \(\Gamma\)
ranges over typing contexts. Primes and numeric subscripts distinguish
metavariables of the same category. Uppercase labels such as \(A\) are a
typographical convention for type members and do not form a separate label
sort.

A member signature is either a proper type or an interval. The two forms
belong to distinct sorts, which the mechanization tracks intrinsically; the
sort index is suppressed in the paper. A term-member definition \(z\) is
paired with a proper signature, whereas a type definition \(T\) is paired
with an interval signature. Labels occupy one syntactic name space.

The binder in a function scopes over its codomain. The binder \(x\) in
\(\pairty{x}{S}{a}{\tau}\) scopes over the member signature \(\tau\); it
denotes the first component of the pair. We write
\(\subst{\tau}{p}{x}\) for capture-avoiding opening. Singleton types
\(\sing{p}\) record term identity. An abstract selection \(p.A\) denotes an
unknown type constrained by the interval stored at \(A\).

The paper uses names, with \(\Gamma(x)=T\) for context lookup, and treats
well-scopedness of the syntax in each judgment as implicit. Substitution and
renaming are capture avoiding; beneath a binder they act on the free variables
and leave the bound variable fixed.

\subsection{Static semantics}
\label{sec:static-semantics}

The static semantics consists of precise path typing
\(\Gamma\vdash p::\tau\), generalized subtyping
\(\Gamma\vdash\tau<:\tau'\), generalized well-formedness
\(\Gamma\vdash\tau\ \wf\), and term typing \(\Gamma\vdash t:T\). Precise
path typing synthesizes either a proper type or an interval and admits no
subsumption. Term typing assigns a singleton type to a path and provides
subsumption separately. The four judgments are proof-relevant inductive
families: their derivations may be inspected by the semantic interpretation.

The matching selection rule opens the member at the receiver's first
projection. A pair can extend a chain of earlier pairs: if its label differs
from the requested label, \ruleref{p-miss} continues at \(p.1.a\). Dynamic
resolution uses the same convention. Function domains are contravariant,
and codomains are compared under the target domain. Dependent pairs are
covariant in their first component and their member; the member comparison
is made under the source first-component type and applies uniformly to
proper and interval members.

Well-formed intervals have consistent bounds. Term paths receive singleton
types, while pair introduction exposes exact information: a stored location
receives its singleton type and a stored type definition receives equal
lower and upper bounds. Subtyping subsequently hides either component. In
\ruleref{t-let}, the result type \(T\) is formed in the outer context and does
not depend on \(x\).

\subsection{Runtime configurations and transitions}
\label{sec:runtime-configurations}
\label{sec:machine}

A store is an immutable sequence of values. A value appended at a fresh
location may refer only to older cells. We write \(\sigma(x)=v\) for lookup,
\(\sigma\mathbin{\cdot}(z\mapsto v)\) for extension at a fresh location
\(z\), and
\(\mathsf{dom}(\sigma)=\{i\mid i<|\sigma|\}\) for the allocated locations.
Machine configurations are well scoped by their store, so every free runtime
variable names an allocated location. Runtime terms therefore use the source
grammar.

Path resolution follows a path to either a location or a stored type. The
machine rules that require a term location demand a result of the form
\(\loc{x}\) directly from this judgment.

The CK machine evaluates the bound computation of a let before its body. A
location returns through a suspended let body without allocation. A syntactic
value is allocated, after which the body may refer to its fresh location. In the
named presentation, allocation opens the body at that location.

The complete source and dynamic semantics are collected in
\cref{fig:source-syntax,fig:runtime-syntax}. Each judgment is shown with all
of its rules.

\clearpage
\begingroup
\setlength{\abovedisplayskip}{5pt plus 1pt minus 1pt}
\setlength{\belowdisplayskip}{5pt plus 1pt minus 1pt}
\setlength{\intextsep}{4pt}
\captionsetup{skip=6pt}
\rulesetup

\begin{figure}[H]
\noindent
\begin{minipage}[t]{\linewidth}
\plateheader{Source syntax}{}
\noindent
\begin{minipage}[t]{0.48\linewidth}
\begin{tabular*}{\linewidth}{@{}r@{\ }l@{\extracolsep{\fill}}r@{}}
\(x,y,z\) && \textbf{Term variable}
\\
\(a,b,A\) && \textbf{Member label}
\\[0.35em]
\(p,q\coloneqq\) && \textbf{Path}
\\
  & \(x\) & variable
\\
  & \(p.1\) & first projection
\\
  & \(p.a\) & member selection
\\[0.35em]
\(S,T,U\coloneqq\) && \textbf{Proper type}
\\
  & \(\top\) & top
\\
  & \(\bot\) & bottom
\\
  & \((x:S)\to T\) & dependent function
\\
  & \(\pairty{x}{S}{a}{\tau}\) & dependent pair
\\
  & \(\sing{p}\) & singleton
\\
  & \(p.A\) & type selection
\\[0.35em]
\(\tau,\tau'\coloneqq\) && \textbf{Member signature}
\\
  & \(T\) & proper type
\\
  & \(S..T\) & interval
\end{tabular*}
\end{minipage}\hfill
\begin{minipage}[t]{0.48\linewidth}
\begin{tabular*}{\linewidth}{@{}r@{\ }l@{\extracolsep{\fill}}r@{}}
\(\delta\coloneqq\) && \textbf{Member definition}
\\
  & \(z\) & term member
\\
  & \(T\) & type member
\\[0.35em]
\(t,u\coloneqq\) && \textbf{Term}
\\
  & \(p\) & path
\\
  & \(\lambda(x:S).t\) & abstraction
\\
  & \(\pairval{y}{a}{\delta}\) & pair
\\
  & \(p\,q\) & application
\\
  & \(\mathbf{let}\ x=t\ \mathbf{in}\ u\) & let
\\[0.35em]
\(v\coloneqq\) && \textbf{Value}
\\
  & \(\lambda(x:S).t\) & abstraction
\\
  & \(\pairval{y}{a}{\delta}\) & pair
\\[0.35em]
\(\Gamma\coloneqq\) && \textbf{Typing context}
\\
  & \(\ctxempty\) & empty
\\
  & \(\Gamma,x:T\) & term binding
\end{tabular*}
\end{minipage}
\end{minipage}\par
\begin{minipage}[t]{\linewidth}
\plateheader{Precise path typing}{\Gamma\vdash p::\tau}
\typicallabel{p-miss}
\noindent
\begin{minipage}[t]{0.47\linewidth}
\infrule[\ruledef{p-var}]
  {\Gamma(x)=T}
  {\Gamma\vdash x::T}

\infrule[\ruledef{p-fst}]
  {\Gamma\vdash p::\pairty{x}{S}{a}{\tau}}
  {\Gamma\vdash p.1::S}
\end{minipage}\hfill
\begin{minipage}[t]{0.51\linewidth}
\infrule[\ruledef{p-sel}]
  {\Gamma\vdash p::\pairty{x}{S}{a}{\tau}}
  {\Gamma\vdash p.a::\subst{\tau}{p.1}{x}}

\infrule[\ruledef{p-miss}]
  {\Gamma\vdash p::\pairty{x}{S}{b}{\tau'}
   \\
   \Gamma\vdash p.1.a::\tau
   \andalso
   a\ne b}
  {\Gamma\vdash p.a::\tau}
\end{minipage}
\end{minipage}

\plateheader{Generalized well-formedness}{\Gamma\vdash\tau\ \wf}
\noindent
\begin{minipage}[t]{0.47\linewidth}
\typicallabel{wf-single}
\infax[\ruledef{wf-bot}]
  {\Gamma\vdash\bot\ \wf}

\infax[\ruledef{wf-top}]
  {\Gamma\vdash\top\ \wf}

\infrule[\ruledef{wf-single}]
  {\Gamma\vdash p::T}
  {\Gamma\vdash\sing{p}\ \wf}

\infrule[\ruledef{wf-select}]
  {\Gamma\vdash p::\pairty{x}{S}{A}{T..U}}
  {\Gamma\vdash p.A\ \wf}
\end{minipage}\hfill
\begin{minipage}[t]{0.51\linewidth}
\typicallabel{wf-bounds}
\infrule[\ruledef{wf-fun}]
  {\Gamma\vdash S\ \wf
   \andalso
   \Gamma,x:S\vdash T\ \wf}
  {\Gamma\vdash (x:S)\to T\ \wf}

\infrule[\ruledef{wf-pair}]
  {\Gamma\vdash S\ \wf
   \andalso
   \Gamma,x:S\vdash\tau\ \wf}
  {\Gamma\vdash\pairty{x}{S}{a}{\tau}\ \wf}

\infrule[\ruledef{wf-bounds}]
  {\Gamma\vdash S\ \wf
   \andalso
   \Gamma\vdash T\ \wf
   \\
   \Gamma\vdash S<:T}
  {\Gamma\vdash S..T\ \wf}
\end{minipage}

\caption{Source syntax and static semantics: source syntax, precise path
typing, and generalized well-formedness.}
\Description{The complete source grammar, all four precise path-typing
rules, and all seven generalized well-formedness rules.}
\label{fig:source-syntax}
\label{fig:path-typing}
\label{fig:well-formedness}
\end{figure}

\clearpage
\begin{figure}[H]
\ContinuedFloat
\plateheader{Generalized subtyping}{\Gamma\vdash\tau<:\tau'}
\noindent
\begin{minipage}[t]{0.47\linewidth}
\typicallabel{s-trans}
\infax[\ruledef{s-refl}]
  {\Gamma\vdash\tau<:\tau}

\infrule[\ruledef{s-trans}]
  {\Gamma\vdash\tau_1<:\tau_2
   \andalso
   \Gamma\vdash\tau_2<:\tau_3}
  {\Gamma\vdash\tau_1<:\tau_3}

\infax[\ruledef{s-bot}]
  {\Gamma\vdash\bot<:T}

\infax[\ruledef{s-top}]
  {\Gamma\vdash T<:\top}
\end{minipage}\hfill
\begin{minipage}[t]{0.51\linewidth}
\typicallabel{s-sel-hi}
\infrule[\ruledef{s-widen}]
  {\Gamma\vdash p::T}
  {\Gamma\vdash\sing{p}<:T}

\infrule[\ruledef{s-alias}]
  {\Gamma\vdash p::\sing{q}}
  {\Gamma\vdash\sing{q}<:\sing{p}}

\infrule[\ruledef{s-sel-hi}]
  {\Gamma\vdash p.A::S..T
   \andalso
   \Gamma\vdash S<:T}
  {\Gamma\vdash p.A<:T}

\infrule[\ruledef{s-sel-lo}]
  {\Gamma\vdash p.A::S..T
   \andalso
   \Gamma\vdash S<:T}
  {\Gamma\vdash S<:p.A}
\end{minipage}\par
\begin{minipage}[t]{\linewidth}
\typicallabel{s-bounds}
\infrule[\ruledef{s-fun}]
  {\Gamma\vdash S'<:S
   \andalso
   \Gamma,x:S'\vdash T<:T'}
  {\Gamma\vdash (x:S)\to T<:(x:S')\to T'}

\infrule[\ruledef{s-pair}]
  {\Gamma\vdash S<:S'
   \andalso
   \Gamma,x:S\vdash\tau<:\tau'}
  {\Gamma\vdash
    \pairty{x}{S}{a}{\tau}<:
    \pairty{x}{S'}{a}{\tau'}}

\infrule[\ruledef{s-bounds}]
  {\Gamma\vdash S'<:S
   \andalso
   \Gamma\vdash T<:T'
   \andalso
   \Gamma\vdash S<:T}
  {\Gamma\vdash S..T<:S'..T'}
\end{minipage}

\caption{Source syntax and static semantics (continued): generalized
subtyping.}
\Description{All eleven generalized subtyping rules.}
\label{fig:subtyping}
\end{figure}

\begin{figure}[H]
\plateheader{Term typing}{\Gamma\vdash t:T}
\noindent
\begin{minipage}[t]{0.40\linewidth}
\typicallabel{t-path}
\infrule[\ruledef{t-path}]
  {\Gamma\vdash p::T}
  {\Gamma\vdash p:\sing{p}}

\infrule[\ruledef{t-abs}]
  {\Gamma,x:S\vdash t:T
   \andalso
   \Gamma\vdash S\ \wf}
  {\Gamma\vdash\lambda(x:S).t:(x:S)\to T}
\typicallabel{t-app}

\infrule[\ruledef{t-app}]
  {\Gamma\vdash p:(x:S)\to T
   \andalso
   \Gamma\vdash q:S}
  {\Gamma\vdash p\,q:\subst{T}{q}{x}}
\end{minipage}\hfill
\begin{minipage}[t]{0.57\linewidth}
\typicallabel{t-type-pair}
\infax[\ruledef{t-pair}]
  {\Gamma\vdash
    \pairval{y}{a}{z}:
    \pairty{x}{\sing{y}}{a}{\sing{z}}}

\infrule[\ruledef{t-type-pair}]
  {\Gamma\vdash T\ \wf}
  {\Gamma\vdash
    \pairval{y}{A}{T}:
    \pairty{x}{\sing{y}}{A}{T..T}}
\typicallabel{t-let}
\infrule[\ruledef{t-let}]
  {\Gamma\vdash t:S
   \andalso
   \Gamma\vdash T\ \wf
   \andalso
   \Gamma,x:S\vdash u:T}
  {\Gamma\vdash
    \mathbf{let}\ x=t\ \mathbf{in}\ u:T}

\infrule[\ruledef{t-sub}]
  {\Gamma\vdash t:S
   \andalso
   \Gamma\vdash S<:T
   \andalso
   \Gamma\vdash T\ \wf}
  {\Gamma\vdash t:T}
\end{minipage}

\noindent
\begin{minipage}[t]{0.57\linewidth}
\plateheader{Runtime syntax}{}
\[
\begin{array}{rcl}
r &::=& \loc{x} \mid \typeref{T}
\\
\sigma &::=& \varnothing
  \mid \sigma\mathbin{\cdot}(x\mapsto v)
\\
K &::=& [\,]
  \mid (\mathbf{let}\ x=[\,]\ \mathbf{in}\ u)::K
\\
c &::=& \langle\sigma,K,t\rangle
\end{array}
\]
\[
\referentof{z}=\loc{z}
\qquad
\referentof{T}=\typeref{T}.
\]
\end{minipage}\hfill
\begin{minipage}[t]{0.40\linewidth}
\plateheader{Final state}{\final(c)}
\typicallabel{f-loc}
\infax[\ruledef{f-loc}]
  {\final(\langle\sigma,[\,],x\rangle)}

\infax[\ruledef{f-val}]
  {\final(\langle\sigma,[\,],v\rangle)}
\end{minipage}

\caption{Term typing and run-time configurations: term typing, run-time
syntax, and final configurations.}
\Description{All seven term-typing rules, the complete run-time grammar,
and both final-state rules.}
\label{fig:term-typing}
\label{fig:runtime-syntax}
\end{figure}

\clearpage
\begin{figure}[H]
\ContinuedFloat

\begin{minipage}[t]{\linewidth}
\plateheader{Generalized path resolution}{\sigma\vdash p\resolve r}
\typicallabel{g-miss}
\noindent
\begin{minipage}[t]{0.47\linewidth}
\infax[\ruledef{g-var}]
  {\sigma\vdash x\resolve\loc{x}}

\infrule[\ruledef{g-fst}]
  {\sigma\vdash p\resolve\loc{x}
   \andalso
   \sigma(x)=\pairval{y}{a}{\delta}}
  {\sigma\vdash p.1\resolve\loc{y}}
\end{minipage}\hfill
\begin{minipage}[t]{0.51\linewidth}
\infrule[\ruledef{g-sel}]
  {\sigma\vdash p\resolve\loc{x}
   \andalso
   \sigma(x)=\pairval{y}{a}{\delta}}
  {\sigma\vdash p.a\resolve\referentof{\delta}}

\infrule[\ruledef{g-miss}]
  {\sigma\vdash p\resolve\loc{x}
   \andalso
   \sigma(x)=\pairval{y}{b}{\delta}
   \\
   a\ne b
   \andalso
   \sigma\vdash y.a\resolve r}
  {\sigma\vdash p.a\resolve r}
\end{minipage}
\end{minipage}

\plateheader{Machine transition}
  {\langle\sigma,K,t\rangle\sstep
   \langle\sigma',K',t'\rangle}
\noindent
\begin{minipage}[t]{0.43\linewidth}
\typicallabel{e-app}
\infrule[\ruledef{e-app}]
   {\sigma\vdash p\resolve\loc{x}
   \andalso
   \sigma\vdash q\resolve\loc{y}
   \\
   \sigma(x)=\lambda(z:S).t}
  {\langle\sigma,K,p\,q\rangle
   \sstep\\
   \langle\sigma,K,\subst{t}{y}{z}\rangle}

\infrule[\ruledef{e-path}]
   {\sigma\vdash p\resolve\loc{x}
   \andalso
   p\ \mathsf{not\ a\ variable}}
  {\langle\sigma,K,p\rangle
   \sstep\\
   \langle\sigma,K,x\rangle}
\end{minipage}\hfill
\begin{minipage}[t]{0.54\linewidth}
\typicallabel{e-return}
\infax[\ruledef{e-let}]
  {\langle\sigma,K,\mathbf{let}\ x=t\ \mathbf{in}\ u\rangle
   \sstep\\
   \langle\sigma,
     (\mathbf{let}\ x=[\,]\ \mathbf{in}\ u)::K,t\rangle}

\infax[\ruledef{e-return}]
  {\langle\sigma,
      (\mathbf{let}\ x=[\,]\ \mathbf{in}\ u)::K,y\rangle
   \sstep\\
   \langle\sigma,K,\subst{u}{y}{x}\rangle}
\end{minipage}\par
\begin{minipage}[t]{\linewidth}
\typicallabel{e-alloc}
\infrule[\ruledef{e-alloc}]
  {z\ \mathsf{fresh}}
  {\langle\sigma,
      (\mathbf{let}\ x=[\,]\ \mathbf{in}\ u)::K,v\rangle
   \sstep\\
   \langle
     \sigma\mathbin{\cdot}(z\mapsto v),K,
     \subst{u}{z}{x}\rangle}
\end{minipage}

\caption{Term typing and run-time configurations (continued): generalized
path resolution and CK-machine transitions.}
\Description{All four generalized-resolution rules and all five machine
transitions.}
\label{fig:path-resolution}
\label{fig:machine}
\end{figure}

\endgroup
\clearpage

We write \(\mathsf{initial}(t)=\langle\varnothing,[\,],t\rangle\) for the
initial state of a closed term \(t\).

\section{Proof obligations}
\label{sec:obligations}

\subsection{Term aliases and abstract selections}

The distinction between singleton types and abstract selections is required
by the alias rule. A premise \(\Gamma\vdash p::\sing{q}\) says that \(p\)
has the singleton type of \(q\), and therefore justifies reversing the alias:
\(\sing{q}<:\sing{p}\). A premise \(\Gamma\vdash p::q.A\) only places the
type of \(p\) between the bounds declared by \(q.A\).

An earlier syntax used one type constructor for both roles. Singleton
symmetry could then consume an abstract-membership premise. With
\[
  f=\lambda(x:\top).x,
  \qquad q=\pairval{f}{A}{\top},
  \qquad h=\lambda(z:q.A).q,
\]
the conflated rules admitted the chain
\[
  \sing{q}<:\top<:q.A<:\sing{z}.
\]
They consequently assigned \(\sing{f}\) to \(h\,f\), although the
application returns \(q\). The complete closed program is
\[
\begin{array}{l}
\mathbf{let}\ f=\lambda(x:\top).x\ \mathbf{in}\\
\mathbf{let}\ q=\pairval{f}{A}{\top}\ \mathbf{in}\\
\mathbf{let}\ h=\lambda(z:q.A).q\ \mathbf{in}\\
\mathbf{let}\ g=h\,f\ \mathbf{in}\ g\,f.
\end{array}
\]
Evaluation binds \(g\) to the pair \(q\), while the erroneous singleton
type makes \(g\) usable as the function \(f\). The final application is
stuck because its operator cell contains a pair. Separate constructors prevent an abstract
member assumption from serving as evidence of path identity. All intervals
in the example may have equal bounds, so the bounds-consistency premises do
not address this problem.

\subsection{The open premise of pair subtyping}

The second premise of \ruleref{s-pair} is a derivation under a hypothetical
variable. An assignment of concrete locations to the variables in \(\Gamma\)
contains no location for that variable. Interpreting the premise immediately
would require choosing one before the source pair is observed. Ordinary
substitution lemmas also leave the critical semantic fact unstated: the
chosen location must have type \(S\), and the member referent must have the
source member type opened at the same location.

The run-time realization of a pair contains exactly these facts. If
\(\sigma(z)=\pairval{y}{a}{\delta}\) realizes
\(\pairty{x}{S}{a}{\tau}\), then \(y\) realizes \(S\) and the referent stored
by \(\delta\) realizes \(\tau[y/x]\). The coercion interpretation retains the
member derivation until this evidence becomes available, permitting the
general rule to vary both the first component and the member signature.

\subsection{Proof architecture}

The proof passes from source derivations to store-local evidence in three
stages. First, typed path resolution turns a precise path-typing derivation
into a run-time resolution together with a realization of its result.
Second, interpretation attaches a closed coercion to a subtyping derivation, and
coercion action maps a realization of the source signature to a realization
of the target signature. Premises beneath binders are saved until a concrete
location is available. Function codomains and pair members use distinct saved
evidence because they are consumed at different points in the proof.
Third, interpretation of term typing produces syntax-directed evidence for
the current machine term. That evidence, together with similarly
syntax-directed continuation evidence, is the invariant used for progress and
preservation. Sections~\ref{sec:coercions} and~\ref{sec:invariant} give the
output-annotated rules before using them in the soundness proof.

\section{Store-local semantics}
\label{sec:realization}
\label{sec:coercions}

\subsection{Runtime equality and conversion}

Generalized resolution may end at a value location or at a stored type.
Runtime path equality \(p\equiv_\sigma q\) is the least equivalence that
relates paths resolving to a common referent and is preserved by
capture-avoiding substitution into a path. Runtime conversion
\(\tau\simeq_\sigma\tau'\) is the least sort-preserving equivalence that
permits the analogous substitution into a member signature. More precisely,
if \(p_0\) or \(\tau_0\) is formed with a distinguished free variable \(z\),
the congruence rules replace every free occurrence of \(z\) simultaneously:
\(
  \subst{p_0}{p}{z}\equiv_\sigma\subst{p_0}{q}{z}
\)
and
\(
  \subst{\tau_0}{p}{z}\simeq_\sigma\subst{\tau_0}{q}{z}.
\)
Substitution alpha-renames an intervening function or pair binder when
necessary.
For types beneath an unassigned binder \(z\), let
\(\equiv_{\sigma;z}\) be the least path congruence that contains ambient
run-time equality and relates \(z\) only to itself. Its generators are
\[
  \frac{p\equiv_\sigma q}{p\equiv_{\sigma;z}q}
  \quad(z\notin\mathsf{fv}(p)\cup\mathsf{fv}(q)),
  \qquad
  z\equiv_{\sigma;z}z,
\]
followed by reflexivity, symmetry, transitivity, and path substitution.
Scoped conversion \(\tau\simeq_{\sigma;z}\tau'\) uses the four
run-time-conversion rules in \cref{fig:semantic-judgments}, with
\(\equiv_{\sigma;z}\) in the replacement rule. This definition preserves
ambient conversions beneath the binder without equating the bound variable
with any store path. Deferred function coercions use its proper-type instance;
pair members may also use its interval instance.

\subsection{Realization and coercions}

Fix a source context \(\Gamma\) and a store \(\sigma\). A valuation has the
signature
\[
  \rho:\mathsf{dom}(\Gamma)\longrightarrow\mathsf{dom}(\sigma).
\]
It is a total assignment of the source variables bound by \(\Gamma\) to
allocated store locations; distinct variables may be assigned the same
location. For any path, type, signature, member definition, or term \(X\),
\(\rho(X)\) denotes the capture-avoiding renaming that replaces each free
source variable \(x\) by \(\rho(x)\). The action is homomorphic; for paths,
\[
\rho(p.1)=\rho(p).1,
\qquad
\rho(p.a)=\rho(p).a.
\]
Beneath a function, pair, or let binder, the bound variable is left unchanged
and \(\rho\) acts on the remaining free variables. If
\(z\notin\mathsf{dom}(\Gamma)\) and \(y\in\mathsf{dom}(\sigma)\), valuation
extension has signature
\[
  \rho[z\mapsto y]:
  \mathsf{dom}(\Gamma,z:S)\longrightarrow\mathsf{dom}(\sigma),
\]
with \(\rho[z\mapsto y](z)=y\) and
\(\rho[z\mapsto y](x)=\rho(x)\) for every
\(x\in\mathsf{dom}(\Gamma)\). For any such \(X\) formed beneath the displayed
binder \(z\), \(\rho(X)\) leaves \(z\) bound and the extension satisfies
\[
  \subst{\rho(X)}{y}{z}=(\rho[z\mapsto y])(X).
\]

The soundness proof uses the following store-local judgments:
\[
\begin{array}{lll}
\rho\realizes_\sigma\Gamma
  &\quad&\text{\(\rho\) realizes the source context \(\Gamma\),}\\
\bodycl{\sigma}{S}{t}{T}
  &&\text{\(t\) is a saved body with parameter type \(S\) and result \(T\),}\\
\sigma\realizes x:T
  &&\text{location \(x\) realizes proper type \(T\),}\\
\sigma\realizes r:\tau
  &&\text{referent \(r\) realizes signature \(\tau\),}\\
\sigma\vdash\tau\coerce\tau'
  &&\text{closed coercion from \(\tau\) to \(\tau'\),}\\
\sigma;z:S\vdash T\coerce T'
  &&\text{deferred function-codomain coercion,}\\
\membercl{\sigma}{z}{S}{\tau}{\tau'}
  &&\text{saved pair-member comparison beneath \(z:S\).}
\end{array}
\]
Here \emph{closed} means that every free source variable has been mapped to
a location in \(\sigma\); a closed coercion may therefore mention store
locations. A deferred codomain coercion waits until function application
supplies an argument for \(z\). A pair-member closure waits for the first
component already stored in the pair. The seven families are finite
inductive evidence defined simultaneously. Function and interval
realizations contain coercions, while coercions retain realization evidence
or saved source derivations.

The complete definitions are collected in
\cref{fig:semantic-judgments}.

Runtime path equality preserves the complete resolution graph: if
\(p\equiv_\sigma q\), then
\(\sigma\vdash p\resolve r\) exactly when
\(\sigma\vdash q\resolve r\). Runtime conversion is used when source syntax
names a path such as \(p.1\), while resolution exposes the equal store
location \(y\). It changes semantic evidence only; no machine transition
performs a conversion.

Transporting realization evidence across a runtime conversion uses two
indexed families of structural evidence:
\(\mathsf{SConv}_\sigma(\tau,\tau')\) for signatures and
\(\mathsf{PConv}_\sigma(T,T')\) for proper types. Their scoped forms
\(\mathsf{SConv}_{\sigma;z}\) and \(\mathsf{PConv}_{\sigma;z}\) use the
scoped path relation \(\equiv_{\sigma;z}\). The output judgment
\[
X\Longmapsto_{\mathsf{str}}Y,\qquad
X:\tau\simeq_\sigma\tau',\qquad
Y:\mathsf{SConv}_\sigma(\tau,\tau')
\]
normalizes raw conversion evidence to structural evidence. Its complete
definition appears in \cref{fig:conversion-normalization}.

\begin{lemma}[Normalization and constructor inversion]
\label{lem:conversion-structure}
If \(X:\tau\simeq_\sigma\tau'\), then some \(Y\) satisfies
\(X\Longmapsto_{\mathsf{str}}Y\). Consequently, \(\tau\) and \(\tau'\) have
the same outer constructor. Functions, pairs, and intervals relate their
corresponding components; pair labels agree, and singleton or selection
components contain runtime-equal paths. Beneath a binder, ambient path
equalities are renamed into the extended scope and the bound variable is
related to itself.
\end{lemma}

The replacement case follows the signature syntax; the transitivity case
composes component decompositions.

We write
\[
  X,R\vdash_{\mathsf{conv}}R'
\]
for transport of realization evidence along a raw runtime conversion. The
structural location judgment is written
\(Y_T,R_x\vdash_{\mathsf{cloc}}R_x'\). Their signatures are
\[
\begin{array}{rcl}
X:\tau\simeq_\sigma\tau',\quad
R:\sigma\realizes r:\tau
&\Longrightarrow&
R':\sigma\realizes r:\tau',
\\[2pt]
Y_T:\mathsf{PConv}_\sigma(T,T'),\quad
R_x:\sigma\realizes x:T
&\Longrightarrow&
R_x':\sigma\realizes x:T'.
\end{array}
\]
The first judgment invokes \(X\Longmapsto_{\mathsf{str}}Y\); the second follows the
resulting proper-type structure. Opening a structural component beneath a
pair binder yields an ordinary runtime conversion, so recursive transport of
the stored member returns to the first judgment.

\begin{lemma}[Conversion transport]
\label{lem:conversion-transport}
If \(X:\tau\simeq_\sigma\tau'\) and
\(R:\sigma\realizes r:\tau\), then there is an \(R'\) such that
\(X,R\vdash_{\mathsf{conv}}R'\).
\end{lemma}

The rules in \cref{fig:conversion-transport} first obtain the structural
decomposition supplied by \cref{lem:conversion-structure}, then follow the
realization evidence. Scoped conversion evidence opens at any location:
\[
  X:\tau\simeq_{\sigma;z}\tau'
  \quad\Longrightarrow\quad
  X[y/z]:\subst{\tau}{y}{z}\simeq_\sigma\subst{\tau'}{y}{z}.
\]

For \(E:\rho\realizes_\sigma\Gamma\), brackets retain the displayed
derivation for reannotation under an extended environment; rule labels name
the resulting evidence. The suspended outputs have signatures
\[
\begin{array}{rcl}
E\mid[\Gamma,z:S\vdash t:T]\Longmapsto B
  &\quad& B:\bodycl{\sigma}{\rho(S)}{\rho(t)}{\rho(T)},\\
E\mid[\Gamma,z:S\vdash\tau<:\tau']\Longmapsto M
  &\quad& M:\membercl{\sigma}{z}{\rho(S)}{\rho(\tau)}{\rho(\tau')},\\
E\mid[\Gamma,z:S\vdash T<:T']\Longmapsto F
  &\quad& F:\sigma;z:\rho(S)\vdash\rho(T)\coerce\rho(T').
\end{array}
\]

\clearpage
\begingroup
\setlength{\abovedisplayskip}{5pt plus 1pt minus 1pt}
\setlength{\belowdisplayskip}{5pt plus 1pt minus 1pt}
\setlength{\intextsep}{4pt}
\captionsetup{skip=6pt}
\rulesetup

\begin{figure}[H]
\noindent
\begin{minipage}[t]{0.48\linewidth}
\plateheader{Path equality}{p\equiv_\sigma q}
\typicallabel{re-resolve}
\infax[\ruledef{re-refl}]
  {p\equiv_\sigma p}

\infrule[\ruledef{re-symm}]
  {p\equiv_\sigma q}
  {q\equiv_\sigma p}

\infrule[\ruledef{re-trans}]
  {p\equiv_\sigma q
   \andalso
   q\equiv_\sigma p'}
  {p\equiv_\sigma p'}

\infrule[\ruledef{re-resolve}]
  {\sigma\vdash p\resolve r
   \andalso
   \sigma\vdash q\resolve r}
  {p\equiv_\sigma q}

\infrule[\ruledef{re-cong}]
  {p\equiv_\sigma q}
  {\subst{p_0}{p}{z}\equiv_\sigma\subst{p_0}{q}{z}}
\end{minipage}\hfill
\begin{minipage}[t]{0.50\linewidth}
\plateheader{Runtime conversion}{\tau\simeq_\sigma\tau'}
\typicallabel{rc-replace}
\infax[\ruledef{rc-refl}]
  {\tau\simeq_\sigma\tau}

\infrule[\ruledef{rc-symm}]
  {\tau\simeq_\sigma\tau'}
  {\tau'\simeq_\sigma\tau}

\infrule[\ruledef{rc-trans}]
  {\tau_1\simeq_\sigma\tau_2
   \andalso
   \tau_2\simeq_\sigma\tau_3}
  {\tau_1\simeq_\sigma\tau_3}

\infrule[\ruledef{rc-replace}]
  {p\equiv_\sigma q}
  {\subst{\tau_0}{p}{z}\simeq_\sigma
   \subst{\tau_0}{q}{z}}
\end{minipage}\par
\begin{minipage}[t]{0.48\linewidth}
\plateheader{Environment realization}{\rho\realizes_\sigma\Gamma}
\typicallabel{e-env}
\infrule[\ruledef{e-env}]
  {\forall x\in\mathsf{dom}(\Gamma).\;
   \sigma\realizes\rho(x):\rho(\Gamma(x))}
  {\rho\realizes_\sigma\Gamma}

\end{minipage}\hfill
\begin{minipage}[t]{0.50\linewidth}
\plateheader{Body closure}
  {E\mid[\Gamma,z:S\vdash t:T]\Longmapsto B}
\typicallabel{bc-source}
\infrule[\ruledef{bc-source}]
  {E:\rho\realizes_{\sigma}\Gamma
   \andalso
   \Gamma,z:S\vdash t:T}
  {E\mid[\Gamma,z:S\vdash t:T]
   \Longmapsto\ruleuse{bc-source}}

\end{minipage}

\plateheader{Referent realization}{\sigma\realizes r:\tau}
\noindent
\begin{minipage}[t]{0.42\linewidth}
\typicallabel{e-type}
\infrule[\ruledef{e-loc}]
  {\sigma\realizes x:T}
  {\sigma\realizes\loc x:T}

\end{minipage}\hfill
\begin{minipage}[t]{0.55\linewidth}
\typicallabel{e-type}
\infrule[\ruledef{e-type}]
  {\sigma\vdash S\coerce T
   \andalso
   \sigma\vdash T\coerce U}
  {\sigma\realizes\typeref T:S..U}
\end{minipage}

\plateheader{Location realization}{\sigma\realizes x:T}
\noindent
\begin{minipage}[t]{0.47\linewidth}
\typicallabel{l-single}
\infax[\ruledef{l-top}]
  {\sigma\realizes x:\top}

\infrule[\ruledef{l-fun}]
  {\sigma(x)=\lambda(z:S_0).t
   \andalso
   \bodycl{\sigma}{S_0}{t}{T_0}
   \\
   \sigma\vdash S\coerce S_0
   \andalso
   \sigma;z:S\vdash T_0\coerce T}
  {\sigma\realizes x:(z:S)\to T}
\end{minipage}\hfill
\begin{minipage}[t]{0.50\linewidth}
\typicallabel{l-select}
\infrule[\ruledef{l-pair}]
  {\sigma(x)=\pairval{y}{a}{\delta}
   \andalso
   \sigma\realizes y:S
   \\
   \sigma\realizes\referentof{\delta}:\subst{\tau}{y}{z}}
  {\sigma\realizes x:\pairty{z}{S}{a}{\tau}}

\infrule[\ruledef{l-single}]
  {\sigma\vdash p\resolve\loc x}
  {\sigma\realizes x:\sing p}

\infrule[\ruledef{l-select}]
  {\sigma\vdash p.A\resolve\typeref T
   \andalso
   \sigma\realizes x:T}
  {\sigma\realizes x:p.A}
\end{minipage}
\caption{Store-local judgments (part 1): equality, conversion,
environments, and realization.}
\Description{All five runtime-equality rules, all four runtime-conversion
rules, the environment rule, the function-body rule, both referent rules, and
all five location-realization rules.}
\label{fig:runtime-eq}
\label{fig:semantic-judgments}
\end{figure}

\clearpage
\begin{figure}[H]
\ContinuedFloat
\setlength{\afterruleskip}{1.65em plus 0.15em minus 0.1em}
\plateheader{Pair-member closure}
  {E\mid[\Gamma,z:S\vdash\tau<:\tau']\Longmapsto M}
\typicallabel{mc-source}
\infrule[\ruledef{mc-source}]
  {E:\rho\realizes_\sigma\Gamma
   \andalso
   \Gamma,z:S\vdash\tau<:\tau'}
  {E\mid[\Gamma,z:S\vdash\tau<:\tau']
   \Longmapsto\ruleuse{mc-source}}

\plateheader{Closed semantic coercion}
  {\sigma\vdash \tau_1\coerce \tau_2}
\noindent
\begin{minipage}[t]{0.43\linewidth}
\typicallabel{m-trans}
\infax[\ruledef{m-refl}]
  {\sigma\vdash \tau\coerce \tau}

\infrule[\ruledef{m-trans}]
  {\sigma\vdash \tau_1\coerce \tau_2
   \andalso
   \sigma\vdash \tau_2\coerce \tau_3}
  {\sigma\vdash \tau_1\coerce \tau_3}

\infrule[\ruledef{m-conv}]
  {\tau_1\simeq_\sigma\tau_2}
  {\sigma\vdash \tau_1\coerce \tau_2}

\infax[\ruledef{m-bot}]
  {\sigma\vdash\bot\coerce T}

\infax[\ruledef{m-top}]
  {\sigma\vdash T\coerce\top}
\end{minipage}\hfill
\begin{minipage}[t]{0.54\linewidth}
\typicallabel{m-sel-hi}
\infrule[\ruledef{m-widen}]
  {\sigma\vdash p\resolve\loc x
   \andalso
   \sigma\realizes x:T}
  {\sigma\vdash\sing p\coerce T}

\infrule[\ruledef{m-alias}]
  {\sigma\vdash p\resolve\loc x
   \andalso
   \sigma\vdash q\resolve\loc x}
  {\sigma\vdash\sing q\coerce\sing p}

\infrule[\ruledef{m-sel-lo}]
  {\sigma\vdash p.A\resolve\typeref T
   \andalso
   \sigma\vdash S\coerce T}
  {\sigma\vdash S\coerce p.A}

\infrule[\ruledef{m-sel-hi}]
  {\sigma\vdash p.A\resolve\typeref T
   \andalso
   \sigma\vdash T\coerce U}
  {\sigma\vdash p.A\coerce U}
\end{minipage}\par
\begin{minipage}[t]{\linewidth}
\typicallabel{m-bounds}
\infrule[\ruledef{m-fun}]
  {\sigma\vdash S'\coerce S
   \andalso
   \sigma;z:S'\vdash T\coerce T'}
  {\sigma\vdash
    (z:S)\to T\coerce(z:S')\to T'}

\infrule[\ruledef{m-pair}]
  {\sigma\vdash S\coerce S'
   \andalso
   \membercl{\sigma}{z}{S}{\tau}{\tau'}}
  {\sigma\vdash
    \pairty{z}{S}{a}{\tau}\coerce
    \pairty{z}{S'}{a}{\tau'}}

\infrule[\ruledef{m-bounds}]
  {\sigma\vdash S'\coerce S
   \andalso
   \sigma\vdash T\coerce T'
   \andalso
   \sigma\vdash S\coerce T}
  {\sigma\vdash S..T\coerce S'..T'}
\end{minipage}

\plateheader{Deferred function coercion}
  {\sigma;z:S\vdash T\coerce T'}
\noindent
\begin{minipage}[t]{0.46\linewidth}
\typicallabel{o-trans}
\infax[\ruledef{o-refl}]
  {\sigma;z:S\vdash T\coerce T}

\infrule[\ruledef{o-trans}]
  {\sigma;z:S\vdash T_1\coerce T_2
   \quad
   \sigma;z:S\vdash T_2\coerce T_3}
  {\sigma;z:S\vdash T_1\coerce T_3}

\infrule[\ruledef{o-conv}]
  {T\simeq_{\sigma;z}T'}
  {\sigma;z:S\vdash T\coerce T'}
\end{minipage}\hfill
\begin{minipage}[t]{0.51\linewidth}
\typicallabel{o-narrow}
\infrule[\ruledef{o-narrow}]
  {\sigma\vdash S'\coerce S
   \andalso
   \sigma;z:S\vdash T\coerce T'}
  {\sigma;z:S'\vdash T\coerce T'}

\plateheader{Codomain closure}
  {E\mid[\Gamma,z:S\vdash T<:T']\Longmapsto F}
\infrule[\ruledef{o-source}]
  {E:\rho\realizes_\sigma\Gamma
   \andalso
   \Gamma,z:S\vdash T<:T'}
  {E\mid[\Gamma,z:S\vdash T<:T']
   \Longmapsto\ruleuse{o-source}}
\end{minipage}
\caption{Store-local judgments (part 2): pair-member closures, closed
coercions, and deferred function coercions.}
\Description{The pair-member closure rule, all twelve rules for closed
semantic coercions, and all five rules for deferred function coercions.}
\end{figure}

\endgroup

We abbreviate \(\ruleuse{m-trans}(C_1,C_2)\) by \(C_1;C_2\); the notation is overloaded
for \ruleref{o-trans} on deferred function coercions.

\begin{figure}[H]
\begingroup
\raggedright
\rulesetup
\setlength{\afterruleskip}{1.1em plus 0.15em minus 0.1em}
\setlength{\labelminsep}{0.4em}
\plateheader{Structural conversion evidence}
  {Y:\mathsf{SConv}_\sigma(\tau,\tau')}
Structural evidence has a signature layer and a proper-type layer:
\[
\begin{array}{rcll}
Y &::=& \mathsf{proper}(Y_T)
       \mid \mathsf{interval}(Y_L,Y_U)
  & \text{signature conversion,}
\\
Y_T &::=& \mathsf{top}\mid\mathsf{bot}
       \mid\mathsf{fun}(Y_S,Y_T)
       \mid\mathsf{pair}(Y_S,Y_\tau)
  & \text{proper-type conversion,}
\\
&& \mathsf{single}(X_p)
       \mid\mathsf{selection}_A(X_p). &
\end{array}
\]
Write \(\mathsf{PConv}_\sigma(T,T')\) for the proper-type layer. The two
layers have the following constructor signatures:
\[
\begin{array}{rcl}
\mathsf{proper}&:&
\mathsf{PConv}_\sigma(T,T')\longrightarrow
\mathsf{SConv}_\sigma(T,T'),
\\
\mathsf{interval}&:&
\mathsf{PConv}_\sigma(L,L')\times\mathsf{PConv}_\sigma(U,U')
\longrightarrow\mathsf{SConv}_\sigma(L..U,L'..U'),
\\
\mathsf{top}&:&\mathsf{PConv}_\sigma(\top,\top),
\qquad
\mathsf{bot}:\mathsf{PConv}_\sigma(\bot,\bot),
\\
\mathsf{fun}&:&
\mathsf{PConv}_\sigma(S,S')\times
\mathsf{PConv}_{\sigma;z}(T,T')
\longrightarrow
\mathsf{PConv}_\sigma((z:S)\to T,(z:S')\to T'),
\\
\mathsf{pair}&:&
\mathsf{PConv}_\sigma(S,S')\times
\mathsf{SConv}_{\sigma;z}(\tau,\tau')
\longrightarrow
\mathsf{PConv}_\sigma(
  \pairty{z}{S}{a}{\tau},\pairty{z}{S'}{a}{\tau'}),
\\
\mathsf{single}&:&
(p\equiv_\sigma q)\longrightarrow
\mathsf{PConv}_\sigma(\sing p,\sing q),
\\
\mathsf{selection}_A&:&
(p\equiv_\sigma q)\longrightarrow
\mathsf{PConv}_\sigma(p.A,q.A).
\end{array}
\]

\plateheader{Constructor inversion}{X\Longmapsto_{\mathsf{str}}Y}
For \(X:\tau\simeq_\sigma\tau'\), the result satisfies
\(Y:\mathsf{SConv}_\sigma(\tau,\tau')\).

\noindent
\begin{minipage}[t]{0.47\linewidth}
\infax[\ruledef{cn-refl}]
  {\ruleuse{rc-refl}_\tau\Longmapsto_{\mathsf{str}}\mathsf{id}_\tau}

\infrule[\ruledef{cn-symm}]
  {X\Longmapsto_{\mathsf{str}}Y}
  {\ruleuse{rc-symm}(X)\Longmapsto_{\mathsf{str}}Y^{-1}}
\end{minipage}\hfill
\begin{minipage}[t]{0.51\linewidth}
\infrule[\ruledef{cn-trans}]
  {X_1\Longmapsto_{\mathsf{str}}Y_1
   \andalso
   X_2\Longmapsto_{\mathsf{str}}Y_2}
  {\ruleuse{rc-trans}(X_1,X_2)\Longmapsto_{\mathsf{str}}
   Y_1\mathbin{\cdot}Y_2}

\infax[\ruledef{cn-replace}]
  {\ruleuse{rc-replace}(\tau_0,X_p)
   \Longmapsto_{\mathsf{str}}\mathsf{subst}_{\tau_0}(X_p)}
\end{minipage}

The four structural operations have signatures
\[
\begin{array}{rcl}
\mathsf{id}_\tau
  &:& \mathsf{SConv}_\sigma(\tau,\tau),
\\
(-)^{-1}
  &:& \mathsf{SConv}_\sigma(\tau,\tau')
       \longrightarrow\mathsf{SConv}_\sigma(\tau',\tau),
\\
(-)\mathbin{\cdot}(-)
  &:& \mathsf{SConv}_\sigma(\tau_1,\tau_2)\times
       \mathsf{SConv}_\sigma(\tau_2,\tau_3)
       \longrightarrow\mathsf{SConv}_\sigma(\tau_1,\tau_3),
\\
\mathsf{subst}_{\tau_0}(X_p)
  &:& \mathsf{SConv}_\sigma(
       \subst{\tau_0}{p}{z},\subst{\tau_0}{q}{z})
       \quad(X_p:p\equiv_\sigma q).
\end{array}
\]
Identity, inverse, and composition have the analogous
\(\mathsf{PConv}\) signatures. They are defined structurally over the common
outer constructor. \(\mathsf{subst}\) follows \(\tau_0\), inserts \(X_p\) at every
occurrence of \(z\), uses identity evidence elsewhere, and leaves a nested
binder bound. Both layers have canonical embeddings into raw conversion:
\(\lceil Y\rceil:\tau\simeq_\sigma\tau'\) and
\(\lceil Y_T\rceil:T\simeq_\sigma T'\). The embedding maps identity,
inverse, and composition to \ruleref{rc-refl}, \ruleref{rc-symm}, and
\ruleref{rc-trans}; the remaining constructors use the corresponding derived
congruence of \ruleref{rc-replace}. It is parametric in the path relation, so
at scoped components
\(\lceil Y_T\rceil:T\simeq_{\sigma;z}T'\) and
\(\lceil Y\rceil:\tau\simeq_{\sigma;z}\tau'\).
\endgroup
\caption{Constructor inversion as an output-annotated judgment.}
\Description{The two layers of structural conversion evidence, the four
normalization rules for raw conversion derivations, and the signatures of the
structural operations.}
\label{fig:conversion-normalization}
\end{figure}

\clearpage

\begin{figure}[H]
\begingroup
\raggedright
\rulesetup
\setlength{\afterruleskip}{1.2em plus 0.15em minus 0.1em}
\setlength{\labelminsep}{0.2em}
\captionsetup{skip=6pt}
\plateheader{Transport along runtime conversion}
  {X,R\vdash_{\mathsf{conv}}R'}
For \(X_p:p\equiv_\sigma q\) and
\(Q:\sigma\vdash p\resolve r\), write
\(\mathsf{res}(X_p,Q):\sigma\vdash q\resolve r\) for resolution invariance;
let
\(X_{p.A}:p.A\equiv_\sigma q.A\) be its congruence instance.
Write \(B_\lambda\) for a store-lookup witness
\(\sigma(x)=\lambda(z:S_0).t\), and \(B_p\) for a witness
\(\sigma(x)=\pairval{y}{a}{\delta}\).

Superscript \({}^{-1}\) denotes structural symmetry; \(\lceil-\rceil\)
embeds structural evidence into raw conversion.

\typicallabel{tr-interval}
\infrule[\ruledef{tr-proper}]
  {X\Longmapsto_{\mathsf{str}}\mathsf{proper}(X_T)
   \andalso
   X_T,R_x\vdash_{\mathsf{cloc}}R_x'}
  {X,\ruleuse{e-loc}(R_x)
   \vdash_{\mathsf{conv}}\ruleuse{e-loc}(R_x')}

\infrule[\ruledef{tr-interval}]
  {X\Longmapsto_{\mathsf{str}}\mathsf{interval}(X_L,X_U)}
  {X,\ruleuse{e-type}(C_L,C_U)
   \vdash_{\mathsf{conv}}\ruleuse{e-type}(C_L',C_U')}

\plateheader{Structural transport at locations}
  {Y_T,R_x\vdash_{\mathsf{cloc}}R_x'}

\infax[\ruledef{tr-top}]
  {\mathsf{top},\ruleuse{l-top}
   \vdash_{\mathsf{cloc}}\ruleuse{l-top}}

\infax[\ruledef{tr-single}]
  {\mathsf{single}(X_p),\ruleuse{l-single}(Q)
   \vdash_{\mathsf{cloc}}\ruleuse{l-single}(Q')}

\infax[\ruledef{tr-selection}]
  {\mathsf{selection}_A(X_p),\ruleuse{l-select}(Q,R_W)
   \vdash_{\mathsf{cloc}}\ruleuse{l-select}(Q_A,R_W)}

\infax[\ruledef{tr-fun}]
  {\mathsf{fun}(X_S,X_T),
   \ruleuse{l-fun}(B_\lambda,B,C_I,F_O)
   \vdash_{\mathsf{cloc}}
   \ruleuse{l-fun}(B_\lambda,B,C_X;C_I,F')}

\infrule[\ruledef{tr-pair}]
  {X_S,R_y\vdash_{\mathsf{cloc}}R_y'
   \andalso
   \lceil X_\tau\rceil[y/z],R_\delta
     \vdash_{\mathsf{conv}}R_\delta'}
  {\mathsf{pair}(X_S,X_\tau),
   \ruleuse{l-pair}(B_p,R_y,R_\delta)
   \vdash_{\mathsf{cloc}}
   \ruleuse{l-pair}(B_p,R_y',R_\delta')}

\noindent where
\[
\begin{aligned}
C_L'&=\ruleuse{m-conv}(\lceil X_L^{-1}\rceil);C_L,
& C_U'&=C_U;\ruleuse{m-conv}(\lceil X_U\rceil),
\\
Q'&=\mathsf{res}(X_p,Q),
& Q_A&=\mathsf{res}(X_{p.A},Q),
\\
C_X&=\ruleuse{m-conv}(\lceil X_S^{-1}\rceil),
& F'&=\ruleuse{o-narrow}(C_X,F_O);
        \ruleuse{o-conv}(\lceil X_T\rceil).
\end{aligned}
\]

There is no rule for \(\mathsf{bot}\), since \(\bot\) has no location
realization. The function rule reverses the domain conversion before its two
uses. The pair rule opens the member conversion at the stored first component
\(y\). The store-lookup witnesses are unchanged.
\endgroup
\caption{Transport of referent and location realizations along runtime
conversion.}
\Description{Raw conversion is structurally decomposed before transport of
proper and interval referents; location transport covers top, function, pair,
singleton, and selection realizations. Bottom has no realization constructor.}
\label{fig:conversion-transport}
\end{figure}

Rule \ruleref{e-env} makes environment realization pointwise. Thus
\(\rho[z\mapsto y]\realizes_\sigma\Gamma,z:S\) whenever
\(\rho\realizes_\sigma\Gamma\) and
\(\sigma\realizes y:\rho(S)\). Every location names a store cell, so
\ruleref{l-top} needs no separate lookup premise, and \(\bot\) has no
realization rule. Rule \ruleref{l-fun} records the stored abstraction and its
body, a coercion from the advertised argument type to the stored domain, and
a deferred coercion from the stored body type to the advertised result type.
Rule \ruleref{l-pair} records realization of the first component and of the
member referent opened at that component.

A closed coercion contains the evidence needed to preserve realization at the
same referent. Identity and transitivity give identity and composition;
runtime conversion transports realizations along run-time path equality.
Singleton and selection coercions retain concrete resolution evidence.
Interval coercions retain the comparisons needed to reconstruct the target
bounds.

The dependent cases save their binder premises. A deferred function coercion
may be composed, transported by conversion, or narrowed before an application
supplies its argument location. A pair-member closure has only the source form
\ruleref{mc-source}: it records the member derivation and its environment until
action on a pair exposes the stored first component. This distinction
determines the definition order of the semantic operations below.

\subsection{Evidence-producing derivations}
\label{sec:coercion-construction}

Let \(E:\rho\realizes_\sigma\Gamma\). We decorate the source rules with
the store-local evidence produced by their derivations. The two judgments
needed here are
\[
\begin{array}{rcll}
E\mid\Gamma\vdash p::\tau\Longmapsto\langle r,Q,R\rangle
&\quad&
Q:\sigma\vdash\rho(p)\resolve r,
&R:\sigma\realizes r:\rho(\tau),
\\[2pt]
E\mid\Gamma\vdash\tau<:\tau'\Longmapsto C
&&
C:\sigma\vdash\rho(\tau)\coerce\rho(\tau'),
&
\end{array}
\]
The source judgment following the vertical bar is derived by the same rule
as in \cref{fig:path-typing,fig:subtyping}; the right-hand side records the
rule's output. Erasing \(E\mid\), that output, and evidence-only premises
recovers the original derivation.
Evidence depends on the chosen derivation, so a single source judgment may
admit several annotations. The derivation objects remain implicit. Function
codomains and pair members are retained with \(E\) as closures and interpreted
once a location for their bound variable is available.

For \(R_y:\sigma\realizes y:\rho(S)\), the notation
\(E[z\mapsto(y,R_y)]\) denotes evidence
\(\rho[z\mapsto y]\realizes_\sigma\Gamma,z:S\).  Write \(E(x)\) for the
realization supplied by \ruleref{e-env}, and write \(X[y/z]\) for opening
scoped conversion evidence at \(y\).

\begin{figure}[H]
\begingroup
\raggedright
\hfuzz=2pt
\rulesetup
\setlength{\labelminsep}{0.2em}
\plateheader{Precise path typing with evidence}
  {E\mid\Gamma\vdash p::\tau\Longmapsto\langle r,Q,R\rangle}

\noindent
\begin{minipage}[t]{0.39\linewidth}
\infrule[\ruleuse{p-var}]
  {\Gamma(x)=T}
  {\begin{gathered}
   E\mid\Gamma\vdash x::T\Longmapsto\\[-1pt]
   \langle\loc{\rho(x)},\ruleuse{g-var},
   \ruleuse{e-loc}(E(x))\rangle
   \end{gathered}}
\end{minipage}\hfill
\begin{minipage}[t]{0.58\linewidth}
\infrule[\ruleuse{p-fst}]
  {\begin{gathered}
   E\mid\Gamma\vdash p::\pairty{z}{S}{a}{\tau}\Longmapsto\\[-1pt]
   \langle\loc x,Q,
   \ruleuse{e-loc}(\ruleuse{l-pair}(B_p,R_y,R_\delta))\rangle
   \end{gathered}}
  {\begin{gathered}
   E\mid\Gamma\vdash p.1::S\Longmapsto\\[-1pt]
   \langle\loc y,\ruleuse{g-fst}(Q,B_p),
   \ruleuse{e-loc}(R_y)\rangle
   \end{gathered}}
\end{minipage}

\infrule[\ruleuse{p-sel}]
  {E\mid\Gamma\vdash p::\pairty{z}{S}{a}{\tau}\Longmapsto
   \langle\loc x,Q,
   \ruleuse{e-loc}(\ruleuse{l-pair}(B_p,R_y,R_\delta))\rangle
   \\
   X_\tau,R_\delta\vdash_{\mathsf{conv}}R_\delta'}
  {E\mid\Gamma\vdash p.a::\subst{\tau}{p.1}{z}\Longmapsto
   \langle r_\delta,Q_s,R_\delta'\rangle}

\infrule[\ruleuse{p-miss}]
  {E\mid\Gamma\vdash p::\pairty{z}{S}{b}{\tau'}\Longmapsto
   \langle\loc x,Q,
   \ruleuse{e-loc}(\ruleuse{l-pair}(B_p,R_y,R_\delta))\rangle
   \\
   E\mid\Gamma\vdash p.1.a::\tau\Longmapsto
   \langle r,Q_1,R_1\rangle
   \andalso N:a\ne b}
  {E\mid\Gamma\vdash p.a::\tau\Longmapsto
   \langle r,Q_m,R_1\rangle}

\noindent where
\[
\begin{aligned}
Q_f&=\ruleuse{g-fst}(Q,B_p),
&X_p&=\ruleuse{re-resolve}(\ruleuse{g-var},Q_f),
\\
X_\tau&=\ruleuse{rc-replace}(\rho(\tau),X_p),
&Q_s&=\ruleuse{g-sel}(Q,B_p),
\\
r_\delta&=\referentof{\delta},
&Q_y&=\mathsf{selCong}(Q_1,Q_f,\ruleuse{g-var}),
\\
Q_m&=\ruleuse{g-miss}(Q,B_p,N,Q_y).&&
\end{aligned}
\]
Resolution congruence has the signature
\[
\begin{gathered}
Q_a:\sigma\vdash p.a\resolve r,
\quad Q_p:\sigma\vdash p\resolve\loc x,
\quad Q_q:\sigma\vdash q\resolve\loc x
\\[-1pt]
\Longrightarrow\quad
\mathsf{selCong}(Q_a,Q_p,Q_q):\sigma\vdash q.a\resolve r.
\end{gathered}
\]
In \ruleref{p-sel}, \(Q_f\) resolves \(\rho(p).1\) to \(y\), so \(X_p\)
relates \(y\) and \(\rho(p).1\). Rule \ruleref{rc-replace} substitutes these
paths for every occurrence of the member binder in \(\rho(\tau)\), producing
\[
X_\tau:
\subst{\rho(\tau)}{y}{z}\simeq_\sigma
\subst{\rho(\tau)}{\rho(p).1}{z}.
\]
In \ruleref{p-miss}, \(N\) is the label inequality and \(Q_y\) rebases the
tail resolution from \(\rho(p).1\) to the stored first component \(y\).
\endgroup
\caption{Evidence annotations for all four precise path-typing rules.}
\Description{Output-annotated rules for variables, first projection,
matching selection, and missed selection.}
\label{fig:typed-resolution-definition}
\end{figure}

\begin{figure}[H]
\begingroup
\raggedright
\rulesetup
\setlength{\afterruleskip}{0.85em plus 0.1em minus 0.05em}
\plateheader{Subtyping with evidence}
  {E\mid\Gamma\vdash\tau<:\tau'\Longmapsto C}
\noindent
\begin{minipage}[t]{0.46\linewidth}
\typicallabel{s-trans}
\infax[\ruleuse{s-refl}]
  {E\mid\Gamma\vdash\tau<:\tau
   \Longmapsto\ruleuse{m-refl}}

\infrule[\ruleuse{s-trans}]
  {E\mid\Gamma\vdash\tau_1<:\tau_2
     \Longmapsto C_1
   \\
   E\mid\Gamma\vdash\tau_2<:\tau_3
     \Longmapsto C_2}
  {\begin{gathered}
   E\mid\Gamma\vdash\tau_1<:\tau_3\\[-1pt]
   \Longmapsto C_1;C_2
   \end{gathered}}

\infax[\ruleuse{s-bot}]
  {E\mid\Gamma\vdash\bot<:T
   \Longmapsto\ruleuse{m-bot}}

\infax[\ruleuse{s-top}]
  {E\mid\Gamma\vdash T<:\top
   \Longmapsto\ruleuse{m-top}}
\end{minipage}\hfill
\begin{minipage}[t]{0.52\linewidth}
\typicallabel{s-sel-hi}
\infrule[\ruleuse{s-widen}]
  {E\mid\Gamma\vdash p::T\Longmapsto
   \langle\loc x,Q,\ruleuse{e-loc}(R_T)\rangle}
  {E\mid\Gamma\vdash\sing{p}<:T\Longmapsto
   \ruleuse{m-widen}(Q,R_T)}

\infrule[\ruleuse{s-alias}]
  {E\mid\Gamma\vdash p::\sing{q}\Longmapsto\\
   \langle\loc x,Q_p,
   \ruleuse{e-loc}(\ruleuse{l-single}(Q_q))\rangle}
  {E\mid\Gamma\vdash\sing{q}<:\sing{p}
   \Longmapsto
   \ruleuse{m-alias}(Q_p,Q_q)}

\infrule[\ruleuse{s-sel-hi}]
  {E\mid\Gamma\vdash p.A::S..T\Longmapsto\\[-1pt]
   \langle\typeref W,Q,\ruleuse{e-type}(C_L,C_U)\rangle
   \\ \Gamma\vdash S<:T}
  {\begin{gathered}
   E\mid\Gamma\vdash p.A<:T\\[-1pt]
   \Longmapsto\ruleuse{m-sel-hi}(Q,C_U)
   \end{gathered}}

\infrule[\ruleuse{s-sel-lo}]
  {E\mid\Gamma\vdash p.A::S..T\Longmapsto\\[-1pt]
   \langle\typeref W,Q,\ruleuse{e-type}(C_L,C_U)\rangle
   \\ \Gamma\vdash S<:T}
  {\begin{gathered}
   E\mid\Gamma\vdash S<:p.A\\[-1pt]
   \Longmapsto\ruleuse{m-sel-lo}(Q,C_L)
   \end{gathered}}
\end{minipage}

\plateheader{Dependent subtyping and member instantiation}{}

\infrule[\ruleuse{s-fun}]
  {E\mid\Gamma\vdash S'<:S\Longmapsto C_D
   \andalso
   E\mid[\Gamma,z:S'\vdash T<:T']\Longmapsto F}
  {E\mid\Gamma\vdash (z:S)\to T<:(z:S')\to T'\Longmapsto
   \ruleuse{m-fun}(C_D,F)}

\infrule[\ruleuse{s-pair}]
  {E\mid\Gamma\vdash S<:S'\Longmapsto C_S
   \andalso
   E\mid[\Gamma,z:S\vdash\tau<:\tau']\Longmapsto M}
  {E\mid\Gamma\vdash\pairty{z}{S}{a}{\tau}<:
                    \pairty{z}{S'}{a}{\tau'}\Longmapsto
   \ruleuse{m-pair}(C_S,M)}

\infrule[\ruleuse{s-bounds}]
  {E\mid\Gamma\vdash S'<:S\Longmapsto C_1
   \andalso
   E\mid\Gamma\vdash T<:T'\Longmapsto C_2
   \andalso
   E\mid\Gamma\vdash S<:T\Longmapsto C_3}
  {E\mid\Gamma\vdash S..T<:S'..T'
   \Longmapsto
   \ruleuse{m-bounds}(C_1,C_2,C_3)}

The unannotated premise in each selection rule proves consistency of the
source bounds. Its realized interval already contains the lower and upper
coercions used in the output. Square brackets mark a premise retained with
\(E\). Thus \ruleref{s-fun} produces the deferred coercion \(F\) by
\ruleref{o-source}, while \ruleref{s-pair} produces the member closure
\(M\) by \ruleref{mc-source}. Both rules interpret the first premise
immediately and retain the exact dependent premise for later use.

\plateheader{Instantiation of a pair-member closure}
  {M,R_y\vdash_{\mathsf{mem}}C}
For \(M:\membercl{\sigma}{z}{S}{\tau}{\tau'}\) and
\(R_y:\sigma\realizes y:S\), the result satisfies
\(C:\sigma\vdash\subst{\tau}{y}{z}\coerce
\subst{\tau'}{y}{z}\).  The judgment has one rule:
\infrule[\ruleuse{mc-source}]
  {E\mid[\Gamma,z:S\vdash\tau<:\tau']\Longmapsto M
   \andalso
   E[z\mapsto(y,R_y)]\mid\Gamma,z:S\vdash\tau<:\tau'
   \Longmapsto C}
  {M,R_y\vdash_{\mathsf{mem}}C}
\endgroup
\caption{Evidence annotations for all eleven subtyping rules and pair-member
instantiation.}
\Description{Output-annotated subtyping rules, followed by the sole
member-instantiation rule.}
\label{fig:coercion-compilation}
\end{figure}

\subsection{Coercion action}
\label{sec:coercion-action}

The judgment \(C,R\vdash_{\mathsf{act}}R'\) annotates a coercion derivation
with the target realization evidence it constructs from \(R\).

\begin{figure}[H]
\begingroup
\raggedright
\rulesetup
\setlength{\labelminsep}{0.2em}
\plateheader{Action of closed coercion evidence}
  {C,R\vdash_{\mathsf{act}}R'}
For \(C:\sigma\vdash\tau\coerce\tau'\) and
\(R:\sigma\realizes r:\tau\), the result satisfies
\(R':\sigma\realizes r:\tau'\).

\noindent
\begin{minipage}[t]{0.47\linewidth}
\typicallabel{m-trans}
\infax[\ruleuse{m-refl}]
  {\ruleuse{m-refl},R\vdash_{\mathsf{act}}R}

\infrule[\ruleuse{m-trans}]
  {C_1,R\vdash_{\mathsf{act}}R_1
   \andalso
   C_2,R_1\vdash_{\mathsf{act}}R_2}
  {\ruleuse{m-trans}(C_1,C_2),R\vdash_{\mathsf{act}}R_2}

\infrule[\ruleuse{m-conv}]
  {X,R\vdash_{\mathsf{conv}}R'}
  {\ruleuse{m-conv}(X),R\vdash_{\mathsf{act}}R'}

\infax[\ruleuse{m-top}]
  {\ruleuse{m-top},\ruleuse{e-loc}(R_x)
   \vdash_{\mathsf{act}}\ruleuse{e-loc}(\ruleuse{l-top})}
\end{minipage}\hfill
\begin{minipage}[t]{0.51\linewidth}
\typicallabel{m-sel-hi}
\infax[\ruleuse{m-widen}]
  {\ruleuse{m-widen}(Q_{\mathsf w},R_T),R_{\mathsf w}
   \vdash_{\mathsf{act}}\ruleuse{e-loc}(R_T)}

\infax[\ruleuse{m-alias}]
  {\ruleuse{m-alias}(Q_p,Q_q),R_{\mathsf a}
   \vdash_{\mathsf{act}}R_{\mathsf p}}

\infrule[\ruleuse{m-sel-lo}]
  {C_L,\ruleuse{e-loc}(R_x)
   \vdash_{\mathsf{act}}\ruleuse{e-loc}(R_W)}
  {\begin{gathered}
   \ruleuse{m-sel-lo}(Q_\ell,C_L),\ruleuse{e-loc}(R_x)\\[-1pt]
   \vdash_{\mathsf{act}}
   \ruleuse{e-loc}(\ruleuse{l-select}(Q_\ell,R_W))
   \end{gathered}}

\infrule[\ruleuse{m-sel-hi}]
  {C_U,\ruleuse{e-loc}(R_W)
   \vdash_{\mathsf{act}}\ruleuse{e-loc}(R_U)}
  {\ruleuse{m-sel-hi}(Q_u,C_U),R_{\mathsf u}
   \vdash_{\mathsf{act}}\ruleuse{e-loc}(R_U)}
\end{minipage}

\infax[\ruleuse{m-fun}]
  {\ruleuse{m-fun}(C_D,F_C),
   \ruleuse{e-loc}(\ruleuse{l-fun}(B_\lambda,B,C_I,F_O))
   \vdash_{\mathsf{act}}
   \ruleuse{e-loc}\bigl(
     \ruleuse{l-fun}(B_\lambda,B,C_D;C_I,F_O')\bigr)}

\infax[\ruleuse{m-bounds}]
  {\ruleuse{m-bounds}(C_L',C_U',C_0),
   \ruleuse{e-type}(C_L,C_U)
   \vdash_{\mathsf{act}}
   \ruleuse{e-type}(C_L';C_L,C_U;C_U')}

\noindent where
\[
\begin{aligned}
Q_{\mathsf w}&:\sigma\vdash p\resolve\loc x,
& Q_{\mathsf w}'&:\sigma\vdash p\resolve\loc{x_{\mathsf w}},
\\
Q_p&:\sigma\vdash p\resolve\loc x,
& Q_q&:\sigma\vdash q\resolve\loc x,
\\
Q_q'&:\sigma\vdash q\resolve\loc{x_{\mathsf a}},
&
\\
Q_u&:\sigma\vdash p.A\resolve\typeref W,
& Q_u'&:\sigma\vdash p.A\resolve\typeref{W_u},
\\
R_{\mathsf w}&=\ruleuse{e-loc}(\ruleuse{l-single}(Q_{\mathsf w}')),
&
R_{\mathsf a}&=\ruleuse{e-loc}(\ruleuse{l-single}(Q_q')),
\\
R_{\mathsf p}&=\ruleuse{e-loc}(\ruleuse{l-single}(Q_p)),
&
R_{\mathsf u}&=\ruleuse{e-loc}(\ruleuse{l-select}(Q_u',R_W)),
\\
F_O'&=\ruleuse{o-narrow}(C_D,F_O);F_C.&&
\end{aligned}
\]

There is no \ruleref{m-bot} action rule, since no realization of \(\bot\)
can be formed. Resolution determinism gives
\(x_{\mathsf w}=x\), \(x_{\mathsf a}=x\), and \(W_u=W\), identifying the
referents in the widening, aliasing, and upper-selection rules. The function rule composes its domain
coercion with the stored input evidence and retains the codomain coercion in
deferred form. In the bounds rule, \(C_0\) justified construction of the
source interval and is not needed in the result.
\endgroup
\caption{Action of closed coercion evidence (part 1): direct,
path-dependent, function, and interval rules.}
\Description{Output rules for reflexivity, transitivity, conversion, top,
widening, aliasing, both selection coercions, functions, and intervals.}
\label{fig:coercion-operations}
\end{figure}

\clearpage
\begin{figure}[H]
\ContinuedFloat
\begingroup
\raggedright
\rulesetup
\plateheader{Action of closed coercion evidence (continued)}
  {C,R\vdash_{\mathsf{act}}R'}

\infrule[\ruleuse{m-pair}]
  {C_S,\ruleuse{e-loc}(R_y)
     \vdash_{\mathsf{act}}\ruleuse{e-loc}(R_y')
   \andalso
   M,R_y\vdash_{\mathsf{mem}}C_M
   \andalso
   C_M,R_\delta\vdash_{\mathsf{act}}R_\delta'}
  {\ruleuse{m-pair}(C_S,M),
   \ruleuse{e-loc}(\ruleuse{l-pair}(B_p,R_y,R_\delta))
   \vdash_{\mathsf{act}}
   \ruleuse{e-loc}(\ruleuse{l-pair}(B_p,R_y',R_\delta'))}

The pair rule uses the original \(R_y:S\) to instantiate \(M\), while
\(R_y':S'\) reconstructs the first component of the target pair.

\plateheader{Instantiation of deferred function evidence}
  {F,R_y\vdash_{\mathsf{cod}}C}
For \(F:\sigma;z:S\vdash T\coerce T'\) and
\(R_y:\sigma\realizes y:S\), the result satisfies
\(C:\sigma\vdash\subst{T}{y}{z}\coerce\subst{T'}{y}{z}\).

\noindent
\begin{minipage}[t]{0.47\linewidth}
\infax[\ruleuse{o-refl}]
  {\ruleuse{o-refl},R_y\vdash_{\mathsf{cod}}\ruleuse{m-refl}}

\infax[\ruleuse{o-conv}]
  {\ruleuse{o-conv}(X),R_y\vdash_{\mathsf{cod}}
   \ruleuse{m-conv}(X[y/z])}
\end{minipage}\hfill
\begin{minipage}[t]{0.51\linewidth}
\infrule[\ruleuse{o-trans}]
  {F_1,R_y\vdash_{\mathsf{cod}}C_1
   \andalso
   F_2,R_y\vdash_{\mathsf{cod}}C_2}
  {\ruleuse{o-trans}(F_1,F_2),R_y\vdash_{\mathsf{cod}}C_1;C_2}
\end{minipage}

\infrule[\ruleuse{o-narrow}]
  {C_D,\ruleuse{e-loc}(R_y)
     \vdash_{\mathsf{act}}\ruleuse{e-loc}(R_y')
   \andalso
   F,R_y'\vdash_{\mathsf{cod}}C_F}
  {\ruleuse{o-narrow}(C_D,F),R_y\vdash_{\mathsf{cod}}C_F}

\infrule[\ruleuse{o-source}]
  {E\mid[\Gamma,z:S\vdash T<:T']\Longmapsto F
   \andalso
   E[z\mapsto(y,R_y)]\mid\Gamma,z:S\vdash T<:T'
     \Longmapsto C}
  {F,R_y\vdash_{\mathsf{cod}}C}
\endgroup
\caption{Action of closed coercion evidence (part 2): pairs and
instantiation of deferred function evidence.}
\Description{The pair action rule and all five deferred-function rules.}
\end{figure}

Coercion action preserves the referent and changes only its realization
evidence; the machine state and store remain fixed. Pair action instantiates
the saved member derivation after the source pair realization exposes its
stored first component. Deferred function evidence is instantiated when
application supplies its argument location.

\subsection{Totality and semantic subtyping}
\label{sec:coercion-totality}

The rules expose the dependency order directly. Normalization produces
structural conversion evidence; structural location transport follows that
evidence, and raw conversion transport combines the two. Path interpretation
uses transport, and subtyping interpretation uses path interpretation. Member
instantiation resumes subtyping interpretation in an extended environment.
Coercion action may invoke member instantiation in its pair rule.
Deferred-function instantiation is last: its source rule resumes subtyping
interpretation and its narrowing rule invokes coercion action. Function
action itself only rearranges deferred evidence.

\begin{lemma}[Typed path resolution]
\label{lem:typed-resolution}
If \(E:\rho\realizes_\sigma\Gamma\) and
\(\Gamma\vdash p::\tau\), then there exist \(r,Q,R\) such that
\(E\mid\Gamma\vdash p::\tau\Longmapsto\langle r,Q,R\rangle\).
\end{lemma}

The selection cases are the substantive steps: matching selection transports
the stored member realization to the receiver's first projection, while
missed selection rebases the tail resolution at the same first component.

\begin{lemma}[Coercion construction]
\label{lem:coercion-construction}
If \(E:\rho\realizes_\sigma\Gamma\) and
\(\Gamma\vdash\tau_1<:\tau_2\), then there is a coercion \(C\) such
that \(E\mid\Gamma\vdash\tau_1<:\tau_2\Longmapsto C\).
\end{lemma}

The proof is induction on the source derivation using the rules in
\cref{fig:coercion-compilation}.  Typed path resolution supplies the
referent and its precise realization in the widening, alias, and selection
cases.  The dependent premise of a function or pair is stored, rather than
interpreted under an environment that lacks its bound variable.

\begin{lemma}[Pair-member instantiation]
\label{lem:member-instantiation}
If \(R_y:\sigma\realizes y:S\) and
\(M:\membercl{\sigma}{z}{S}{\tau}{\tau'}\), then there is a \(C_M\) such
that \(M,R_y\vdash_{\mathsf{mem}}C_M\).
\end{lemma}

\begin{lemma}[Coercion action]
\label{lem:coercion-action}
If \(C:\sigma\vdash\tau_1\coerce\tau_2\) and
\(R:\sigma\realizes r:\tau_1\), then there is an \(R'\) such that
\(C,R\vdash_{\mathsf{act}}R'\).
\end{lemma}

The referent remains \(r\).  The action judgment changes only the semantic
evidence attached to an existing location or stored type.  Its totality is
the content of the lemma; the well-founded induction needed for the pair rule
is described in \cref{sec:coercion-termination}.

\begin{lemma}[Deferred-codomain instantiation]
\label{lem:codomain-instantiation}
If \(R_y:\sigma\realizes y:S\) and
\(F:\sigma;z:S\vdash T\coerce T'\), then there is a \(C_F\) such that
\(F,R_y\vdash_{\mathsf{cod}}C_F\).
\end{lemma}

Subtyping interpretation followed by coercion action gives the semantic
content of a source subtyping derivation.

\begin{theorem}[Semantic subtyping]
\label{thm:semantic-subtyping}
If \(E:\rho\realizes_\sigma\Gamma\),
\(\Gamma\vdash\tau<:\tau'\), and
\(R:\sigma\realizes r:\rho(\tau)\), then there exist \(C,R'\) such that
\[
  E\mid\Gamma\vdash\tau<:\tau'\Longmapsto C,
  \qquad
  C,R\vdash_{\mathsf{act}}R',
  \qquad
  R':\sigma\realizes r:\rho(\tau').
\]
\end{theorem}

\subsection{Action for the general pair rule}
\label{sec:coercion-pair}

Fix
\[
  E:\rho\realizes_\sigma\Gamma,
  \qquad
  \Gamma\vdash S<:S',
  \qquad
  \Gamma,z:S\vdash\tau<:\tau'.
\]
The output-annotated \ruleref{s-pair} rule gives \(C_S\) and \(M\) with
\[
  E\mid\Gamma\vdash S<:S'\Longmapsto C_S,
  \qquad
  E\mid[\Gamma,z:S\vdash\tau<:\tau']\Longmapsto M,
\]
and hence
\[
  E\mid\Gamma\vdash
    \pairty{z}{S}{a}{\tau}<:\pairty{z}{S'}{a}{\tau'}
  \Longmapsto\ruleuse{m-pair}(C_S,M).
\]
Suppose the source realization exposes
\[
  \sigma(x)=\pairval{y}{a}{\delta},
  \qquad
  R_y:\sigma\realizes y:\rho(S),
  \qquad
  R_\delta:\sigma\realizes\referentof{\delta}:
    \subst{\rho(\tau)}{y}{z}.
\]
The \ruleref{m-pair} action rule has two independent uses of \(R_y\).
Choose \(R_y'\) and \(C_\tau\) by
\[
  C_S,\ruleuse{e-loc}(R_y)
    \vdash_{\mathsf{act}}\ruleuse{e-loc}(R_y'),
  \qquad
  M,R_y\vdash_{\mathsf{mem}}C_\tau.
\]
The first judgment reconstructs the target first component at
\(\rho(S')\).  The second extends \(E\) with the original
\(R_y:\rho(S)\), resumes the suspended member derivation, and yields
\[
  C_\tau:\sigma\vdash
  \subst{\rho(\tau)}{y}{z}\coerce
  \subst{\rho(\tau')}{y}{z}.
\]
Finally choose \(R_\delta'\) with
\(C_\tau,R_\delta\vdash_{\mathsf{act}}R_\delta'\).  The three judgments
are precisely the premises of the pair-action rule.  The argument applies
to proper and interval members; \(\referentof{\delta}\) is respectively a
location or a stored type.

\subsection{Well-foundedness of coercion action}
\label{sec:coercion-termination}

Transitivity and selection recurse on structural subcoercions. Pair action
also invokes the coercion obtained by instantiating its member closure. The
result can be larger than the enclosing pair coercion, so coercion structure
alone does not justify recursion. The allocation order of the immutable store
provides the required decrease.

Give a location referent \(\loc{x}\) its position in the allocation order,
starting at one, and give every type referent rank zero. Write this number as
\(\mathsf{rank}_\sigma(r)\). If
\[
  \sigma(x)=\pairval{y}{a}{\delta},
\]
then both referents stored in the pair have smaller rank:
\[
  \mathsf{rank}_\sigma(\loc y)
    <\mathsf{rank}_\sigma(\loc x),
  \qquad
  \mathsf{rank}_\sigma(\referentof{\delta})
    <\mathsf{rank}_\sigma(\loc x).
\]
For an interval member, \(\referentof{\delta}=\typeref{T}\) has rank zero. For
a proper member it is an older location named by the stored pair.
Fresh allocation preserves the relative rank of all existing referents.

Let \(|C|\) be the number of constructors in a closed coercion; saved
resolution witnesses, member closures, and deferred function coercions do not
contribute to this number.
Existence of an action derivation is proved by well-founded induction on
\[
  \bigl(\mathsf{rank}_\sigma(r),|C|\bigr)
\]
in lexicographic order.  The premises of \ruleref{m-trans},
\ruleref{m-sel-lo}, and \ruleref{m-sel-hi} retain the referent and use a
proper subcoercion, decreasing the secondary component.  In
\ruleref{m-pair}, the first action premise concerns \(\loc y\), and the
second concerns \(\referentof{\delta}\).  Both strictly decrease the primary
component, independently of the size of the coercion obtained by member
instantiation.  These are all recursive action premises.

\section{The machine invariant}
\label{sec:invariant}

\subsection{Syntax-directed run-time evidence}

Source typing permits subsumption at any depth. Direct inversion of that
judgment therefore exposes an arbitrary final \ruleref{t-sub} derivation. The
machine invariant uses syntax-directed evidence whose precise constructors
retain final subsumption as a trailing semantic coercion. Each value
constructor carries such a coercion; path, application, and let term evidence
do likewise, while value terms refer to value evidence. We write
\(\sigma\vdash_{V}v:T\) for values and
\(\sigma\vdash_{R}t:T\) for run-time terms. These judgments are evidence in
the metatheory; they do not extend the machine syntax.

\Cref{fig:term-evidence} gives all constructors. The source types in the two
pair-value rules are their precise introduction types. In
\ruleref{rt-let}, \(U\) is independent of the let binder. Continuations are
typed as transformations from the type of the current term to the final
result type. The empty-continuation rule permits a coercion, absorbing
subsumption at the boundary of the machine stack.

\begin{figure}[H]
\begingroup
\raggedright
\rulesetup
\judgheader{Value evidence}{\sigma\vdash_V v:T}
\noindent
\begin{minipage}[t]{\linewidth}
\typicallabel{rv-abs}
\infrule[\ruledef{rv-abs}]
  {\bodycl{\sigma}{S_0}{t}{T_0}
   \andalso
   \sigma\vdash (x:S_0)\to T_0\coerce T}
  {\sigma\vdash_V\lambda(x:S_0).t:T}
\end{minipage}\par
\begin{minipage}[t]{0.48\linewidth}
\typicallabel{rv-pair}
\infrule[\ruledef{rv-pair}]
  {\sigma\vdash
    \pairty{x}{\sing y}{a}{\sing z}\coerce T}
  {\sigma\vdash_V\pairval{y}{a}{z}:T}
\end{minipage}\hfill
\begin{minipage}[t]{0.50\linewidth}
\typicallabel{rv-type-pair}
\infrule[\ruledef{rv-type-pair}]
  {\sigma\vdash
    \pairty{x}{\sing y}{A}{S..S}\coerce T}
  {\sigma\vdash_V\pairval{y}{A}{S}:T}
\end{minipage}

\judgheader{Term evidence}{\sigma\vdash_R t:T}
\noindent
\begin{minipage}[t]{0.43\linewidth}
\typicallabel{rt-value}
\infrule[\ruledef{rt-path}]
  {\sigma\vdash p\resolve\loc x
   \andalso
   \sigma\vdash\sing p\coerce T}
  {\sigma\vdash_R p:T}

\infrule[\ruledef{rt-value}]
  {\sigma\vdash_V v:T}
  {\sigma\vdash_R v:T}
\end{minipage}\hfill
\begin{minipage}[t]{0.55\linewidth}
\typicallabel{rt-let}
\infrule[\ruledef{rt-app}]
  {\sigma\vdash_R p:(x:S)\to U
   \andalso
   \sigma\vdash_R q:S
   \\
   \sigma\vdash\subst{U}{q}{x}\coerce T}
  {\sigma\vdash_R p\,q:T}

\infrule[\ruledef{rt-let}]
  {\sigma\vdash_R t:S
   \andalso
   \bodycl{\sigma}{S}{u}{U}
   \\
   \sigma\vdash U\coerce T}
  {\sigma\vdash_R\mathbf{let}\ x=t\ \mathbf{in}\ u:T}
\end{minipage}

\judgheader{Continuation evidence}{\sigma\vdash K:S\Rightarrow T}
\noindent
\begin{minipage}[t]{0.43\linewidth}
\typicallabel{rk-hole}
\infrule[\ruledef{rk-hole}]
  {\sigma\vdash S\coerce T}
  {\sigma\vdash[\,]:S\Rightarrow T}
\end{minipage}\hfill
\begin{minipage}[t]{0.55\linewidth}
\typicallabel{rk-let}
\infrule[\ruledef{rk-let}]
  {\sigma\vdash K:U\Rightarrow T
   \andalso
   \bodycl{\sigma}{S}{u}{U'}
   \\
   \sigma\vdash U'\coerce U}
  {\sigma\vdash
    (\mathbf{let}\ x=[\,]\ \mathbf{in}\ u)::K:S\Rightarrow T}
\end{minipage}

\judgheader{State evidence}{\vdash\langle\sigma,K,t\rangle:T}
\typicallabel{rs-state}
\infrule[\ruledef{rs-state}]
  {\sigma\vdash_R t:S
   \andalso
   \sigma\vdash K:S\Rightarrow T}
  {\vdash\langle\sigma,K,t\rangle:T}
\endgroup
\caption{The complete store-local machine invariant: syntax-directed
evidence for values, terms, continuations, and states.}
\Description{All three value rules, four term rules, two continuation rules,
and the state-evidence rule.}
\label{fig:term-evidence}
\label{fig:state-evidence}
\end{figure}

For \(c=\langle\sigma,K,t\rangle\), write \(\vdash c:T\) for the state
evidence judgment in \cref{fig:state-evidence}.

\subsection{Evidence-producing typing derivations}

The same annotation extends to source typing:
\[
  E\mid\Gamma\vdash t:T\Longmapsto R_t,
  \qquad
  R_t:\sigma\vdash_R\rho(t):\rho(T).
\]
Erasing the annotation and its evidence-only premises recovers the source
typing derivation. A body below a binder is retained as a closure, and subsumption
uses the judgment below to append the coercion produced by subtyping
interpretation to the trailing coercion already present in \(R_t\).

\begin{figure}[H]
\begingroup
\raggedright
\rulesetup
\setlength{\afterruleskip}{0.5em plus 0.1em minus 0.05em}
\setlength{\labelminsep}{0.5em}
\plateheader{Appending a trailing coercion}
  {R,C\vdash_{\mathsf{trail}}R'}

\infax[\ruleuse{rt-path}]
  {\ruleuse{rt-path}(Q,C_1),C_2\vdash_{\mathsf{trail}}
   \ruleuse{rt-path}(Q,C_1;C_2)}

\infrule[\ruleuse{rt-value}]
  {R_v,C\vdash_{\mathsf{trail}}R_v'}
  {\ruleuse{rt-value}(R_v),C\vdash_{\mathsf{trail}}
   \ruleuse{rt-value}(R_v')}

\infax[\ruleuse{rt-app}]
  {\ruleuse{rt-app}(R_1,R_2,C_1),C_2
   \vdash_{\mathsf{trail}}
   \ruleuse{rt-app}(R_1,R_2,C_1;C_2)}

\infax[\ruleuse{rt-let}]
  {\ruleuse{rt-let}(R_1,B,C_1),C_2
   \vdash_{\mathsf{trail}}
   \ruleuse{rt-let}(R_1,B,C_1;C_2)}

\infax[\ruleuse{rv-abs}]
  {\ruleuse{rv-abs}(B,C_1),C_2\vdash_{\mathsf{trail}}
   \ruleuse{rv-abs}(B,C_1;C_2)}

\infax[\ruleuse{rv-pair}]
  {\ruleuse{rv-pair}(C_1),C_2\vdash_{\mathsf{trail}}
   \ruleuse{rv-pair}(C_1;C_2)}

\infax[\ruleuse{rv-type-pair}]
  {\ruleuse{rv-type-pair}(C_1),C_2\vdash_{\mathsf{trail}}
   \ruleuse{rv-type-pair}(C_1;C_2)}

\plateheader{Term typing with evidence}
  {E\mid\Gamma\vdash t:T\Longmapsto R_t}
\endgroup
\begingroup
\raggedright
\rulesetup
\setlength{\afterruleskip}{1.2em plus 0.15em minus 0.1em}
\setlength{\labelminsep}{0pt}
\noindent
\begin{minipage}[t]{0.49\linewidth}
\infrule[\ruleuse{t-path}]
  {E\mid\Gamma\vdash p::T\Longmapsto
   \langle\loc x,Q,\ruleuse{e-loc}(R_x)\rangle}
  {E\mid\Gamma\vdash p:\sing{p}\Longmapsto
   \ruleuse{rt-path}(Q,\ruleuse{m-refl})}

\infrule[\ruleuse{t-abs}]
  {E\mid[\Gamma,z:S\vdash t:T]\Longmapsto B
   \andalso
   \Gamma\vdash S\ \wf}
  {E\mid\Gamma\vdash\lambda(z:S).t:(z:S)\to T
   \Longmapsto R_\lambda}
\end{minipage}\hfill
\begin{minipage}[t]{0.505\linewidth}
\infrule[\ruleuse{t-app}]
  {E\mid\Gamma\vdash p:(z:S)\to T\Longmapsto R_1
   \\
   E\mid\Gamma\vdash q:S\Longmapsto R_2}
  {\begin{gathered}
   E\mid\Gamma\vdash p\,q:\subst{T}{q}{z}\Longmapsto\\[-1pt]
   \ruleuse{rt-app}(R_1,R_2,\ruleuse{m-refl})
   \end{gathered}}

\infax[\ruleuse{t-pair}]
  {E\mid\Gamma\vdash
    \pairval{y}{a}{z}:\pairty{x}{\sing{y}}{a}{\sing{z}}
    \Longmapsto R_p}
\end{minipage}

\infrule[\ruleuse{t-type-pair}]
  {\Gamma\vdash T\ \wf}
  {E\mid\Gamma\vdash
   \pairval{y}{A}{T}:\pairty{x}{\sing{y}}{A}{T..T}\Longmapsto
   \ruleuse{rt-value}(
     \ruleuse{rv-type-pair}(\ruleuse{m-refl}))}

\infrule[\ruleuse{t-let}]
  {E\mid\Gamma\vdash t:S\Longmapsto R_1
   \andalso
   \Gamma\vdash T\ \wf
   \\
   E\mid[\Gamma,z:S\vdash u:T]\Longmapsto B}
  {E\mid\Gamma\vdash\mathbf{let}\ z=t\ \mathbf{in}\ u:T
   \Longmapsto
   \ruleuse{rt-let}(R_1,B,
                    \ruleuse{m-refl})}

\infrule[\ruleuse{t-sub}]
  {E\mid\Gamma\vdash t:S\Longmapsto R_1
   \andalso
   E\mid\Gamma\vdash S<:T\Longmapsto C
   \andalso
   \Gamma\vdash T\ \wf
   \\
   R_1,C\vdash_{\mathsf{trail}}R_2}
  {E\mid\Gamma\vdash t:T\Longmapsto R_2}

\noindent where
\[
R_\lambda=
  \ruleuse{rt-value}(\ruleuse{rv-abs}(
    B,\ruleuse{m-refl})),
\qquad
R_p=\ruleuse{rt-value}(\ruleuse{rv-pair}(\ruleuse{m-refl})).
\]

\plateheader{Application of a body closure}
  {B,R_y\vdash_{\mathsf{body}}R_b}
\infrule[\ruleuse{bc-source}]
  {E\mid[\Gamma,z:S\vdash t:T]\Longmapsto B
   \andalso
   E[z\mapsto(y,R_y)]\mid\Gamma,z:S\vdash t:T
     \Longmapsto R_b}
  {B,R_y\vdash_{\mathsf{body}}R_b}
\endgroup
\caption{Trailing coercions, evidence annotations for all seven source
typing rules, and body-closure application.}
\Description{The complete trailing-coercion judgment, all output-annotated
typing rules, and the sole body-application rule.}
\label{fig:typing-interpretation}
\end{figure}

For \(R:\sigma\vdash_R t:S\) and \(C:\sigma\vdash S\coerce T\), trailing
coercion produces \(R':\sigma\vdash_R t:T\); the definition is overloaded
for value evidence. The unannotated premises of the typing rules establish
source well-formedness and produce no run-time evidence. The abstraction and
let rules retain their body premises as \(B\). Applying
\(B:\bodycl{\sigma}{S}{t}{T}\) to \(R_y:\sigma\realizes y:S\) produces
\(R_b:\sigma\vdash_R\subst{t}{y}{z}:\subst{T}{y}{z}\).

The trailing-coercion judgment is total by case analysis on its term or value
evidence.

\begin{lemma}[Semantic typing]
\label{lem:semantic-typing}
If \(E:\rho\realizes_{\sigma}\Gamma\) and
\(\Gamma\vdash t:T\), then there is an \(R_t\) such that
\(E\mid\Gamma\vdash t:T\Longmapsto R_t\).
\end{lemma}

For the empty context and store, combining the lemma with \ruleref{rk-hole}
at reflexivity gives \(\vdash\mathsf{initial}(t):T\).

\begin{figure}[H]
\begingroup
\small
\renewcommand{\arraystretch}{1.28}
\begin{tabular*}{\linewidth}{@{\extracolsep{\fill}}p{0.34\linewidth}p{0.27\linewidth}p{0.33\linewidth}@{}}
\toprule
Input evidence & Judgment or rule & Result evidence \\
\midrule
\(X:\tau\simeq_\sigma\tau'\),
\(R:\sigma\realizes r:\tau\)
& \(X,R\vdash_{\mathsf{conv}}R'\)
& \(R':\sigma\realizes r:\tau'\)
\\
\(E:\rho\realizes_\sigma\Gamma\),
\(\Gamma\vdash p::\tau\)
& \(E\mid\Gamma\vdash p::\tau
      \Longmapsto\langle r,Q,R\rangle\)
& \(r\), \(Q:\sigma\vdash\rho(p)\resolve r\),
  \(R:\sigma\realizes r:\rho(\tau)\)
\\
\(E:\rho\realizes_\sigma\Gamma\),
\(\Gamma\vdash\tau<:\tau'\)
& \(E\mid\Gamma\vdash\tau<:\tau'\Longmapsto C\)
& \(C:\sigma\vdash\rho(\tau)\coerce\rho(\tau')\)
\\
\(E\mid[\Gamma,z:S\vdash\tau<:\tau']\Longmapsto M\),\newline
\(R_y:\sigma\realizes y:\rho(S)\)
& \(M,R_y\vdash_{\mathsf{mem}}C_M\)
& \(C_M:\sigma\vdash\subst{\rho(\tau)}{y}{z}\coerce
                    \subst{\rho(\tau')}{y}{z}\)
\\
\(C:\sigma\vdash\tau\coerce\tau'\),
\(R:\sigma\realizes r:\tau\)
& \(C,R\vdash_{\mathsf{act}}R'\)
& \(R':\sigma\realizes r:\tau'\)
\\
\(E\mid[\Gamma,z:S\vdash T<:T']\Longmapsto F\),\newline
\(R_y:\sigma\realizes y:\rho(S)\)
& \(F,R_y\vdash_{\mathsf{cod}}C_F\)
& \(C_F:\sigma\vdash\subst{\rho(T)}{y}{z}\coerce
                    \subst{\rho(T')}{y}{z}\)
\\
\(R_t:\sigma\vdash_R t:S\),
\(C:\sigma\vdash S\coerce T\)
& \(R_t,C\vdash_{\mathsf{trail}}R_t'\)
& \(R_t':\sigma\vdash_R t:T\)
\\
\(E:\rho\realizes_\sigma\Gamma\),
\(\Gamma\vdash t:T\)
& \(E\mid\Gamma\vdash t:T\Longmapsto R_t\)
& \(R_t:\sigma\vdash_R\rho(t):\rho(T)\)
\\
\(E\mid[\Gamma,z:S\vdash t:T]\Longmapsto B\),\newline
\(R_y:\sigma\realizes y:\rho(S)\)
& \(B,R_y\vdash_{\mathsf{body}}R_b\)
& \(R_b:\sigma\vdash_R\subst{\rho(t)}{y}{z}:
                    \subst{\rho(T)}{y}{z}\)
\\
\(R_t:\sigma\vdash_R t:S\),
\(R_K:\sigma\vdash K:S\Rightarrow T\)
& \ruleref{rs-state}
& \(\vdash\langle\sigma,K,t\rangle:T\)
\\
\bottomrule
\end{tabular*}
\endgroup
\caption{Flow of conversion and typing evidence into the machine invariant.}
\Description{Inputs and outputs of conversion transport, source-rule
annotation, coercion action, binder instantiation, trailing coercion
composition, body application, and state construction.}
\label{fig:evidence-flow}
\end{figure}

For a closed source term, these constructions feed the safety lemmas in the
following order:
\[
\begin{aligned}
\ctxempty\vdash t:T
&\quad\Longrightarrow\quad
\vdash\mathsf{initial}(t):T,
\\
\vdash\mathsf{initial}(t):T
\quad\land\quad \mathsf{initial}(t)\steps c
&\quad\xRightarrow{\text{preservation}}\quad
\vdash c:T,
\\
\vdash c:T
&\quad\xRightarrow{\text{progress}}\quad
\final(c)\ \text{or}\ \exists c'.\;c\sstep c'.
\end{aligned}
\]

Body and member premises cross a binder by saving \(E\) with their source
derivation. Applying a body closure resumes its saved typing derivation;
member and deferred-codomain instantiation resume their saved subtyping
derivations. The latter supplies the codomain coercion used in application
preservation. Rule \ruleref{rs-state} combines term evidence with continuation
evidence. Progress analyzes this state evidence, while
preservation reconstructs it after each machine transition.

\subsection{Progress}

Term evidence for a path \(p:T\), together with
\(\sigma\vdash p\resolve\loc x\), yields \(\sigma\realizes x:T\): construct
the singleton realization from the resolution and apply the trailing coercion.
This observation supplies canonical forms for applications. If the operator
has a function type, inversion of \ruleref{l-fun} produces an abstraction at
the resolved location.

All remaining progress cases follow the outer term constructor. A
non-variable path resolves to its location; a variable is final under an empty
continuation or returns through a suspended let body. Values are final under
an empty continuation and allocate under a nonempty continuation. A let
expression pushes its body onto the continuation.

\begin{theorem}[Progress]
\label{thm:progress}
If \(\vdash c:T\), then either \(\final(c)\) or there is a state
\(c'\) such that \(c\sstep c'\).
\end{theorem}

\subsection{Preservation}

The named presentation chooses a globally fresh location at allocation, so
types formed before allocation remain well scoped without changing their
notation.

\begin{theorem}[One-step preservation]
\label{thm:preservation}
If \(\vdash c:T\) and \(c\sstep c'\), then \(\vdash c':T\).
\end{theorem}

The application case contains the essential dependent reasoning. Suppose
\(p\) resolves to \(x\) and \(q\) resolves to \(y\). Inverting
\ruleref{rt-app} gives operator evidence at \((z:S)\to U\), argument
evidence at \(S\), and a suffix coercion
\(C_R:\sigma\vdash\subst{U}{q}{z}\coerce T\). Preservation constructs
evidence for the reduct in the following order.
\begin{enumerate}
\item Apply the operator's trailing coercion to the singleton realization of
  \(x\). Inversion of \ruleref{l-fun} then exposes the stored abstraction,
  a body closure \(B:\bodycl{\sigma}{S_0}{t}{U_0}\), an input coercion
  \(C_I:\sigma\vdash S\coerce S_0\), and deferred output evidence
  \(F_O:\sigma;z:S\vdash U_0\coerce U\).
\item The argument evidence and its resolution give
  \(R_y:\sigma\realizes y:S\).  By coercion action, choose \(R_y'\) with
  \[
    C_I,\ruleuse{e-loc}(R_y)
      \vdash_{\mathsf{act}}\ruleuse{e-loc}(R_y'),
    \qquad R_y':\sigma\realizes y:S_0.
  \]
\item By body application, choose \(R_b\) with
  \[
    B,R_y'\vdash_{\mathsf{body}}R_b,
    \qquad
    R_b:\sigma\vdash_R\subst{t}{y}{z}:\subst{U_0}{y}{z}.
  \]
\item Instantiating the deferred output with the original argument
  realization yields a \(C_F\) with
  \[
    F_O,R_y\vdash_{\mathsf{cod}}C_F,
    \qquad
    C_F:\sigma\vdash\subst{U_0}{y}{z}\coerce\subst{U}{y}{z}.
  \]
\item Since \(q\) and the variable \(y\) resolve to the same location,
  run-time conversion gives a coercion \(C_q\) from
  \(\subst{U}{y}{z}\) to \(\subst{U}{q}{z}\).  By trailing-coercion
  totality, choose \(R_b'\) with
  \[
    R_b,C_F;C_q;C_R\vdash_{\mathsf{trail}}R_b',
    \qquad
    R_b':\sigma\vdash_R\subst{t}{y}{z}:T.
  \]
\end{enumerate}
Thus the body is interpreted at the stored domain, while the saved codomain
comparison is instantiated at the argument type advertised by the function
realization.

For a path step, \(p\) resolves to \(x\). The target term has the precise type
\(\sing{x}\). Runtime equality gives a conversion from \(\sing{x}\) to
\(\sing{p}\), which composes with the saved path coercion. Let-push records the
saved body closure in the continuation, and return instantiates it with the
returned location. In the allocation case, inversion exposes the value's
precise introduction evidence and trailing coercion. After adding the fresh
cell, the corresponding location-realization constructor and coercion action
show that the new location realizes the advertised value type. The proof then
weakens the saved continuation, body closure, and suffix coercion along the
store extension, applies the body closure at the fresh location, and composes
the weakened suffix with the resulting term evidence.

\subsection{Finite executions}

Write \(\steps\) for the reflexive-transitive closure of \(\sstep\).
Iterating \cref{thm:preservation} transports initial state evidence through
any finite execution. Applying \cref{thm:progress} at its last state gives the
standard non-stuckness statement.

\begin{theorem}[Finite preservation]
\label{thm:finite-preservation}
If \(\ctxempty\vdash t:T\) and
\(\mathsf{initial}(t)\steps c\), then \(\vdash c:T\).
\end{theorem}

\begin{theorem}[Type safety]
\label{thm:safety}
If \(\ctxempty\vdash t:T\) and
\(\mathsf{initial}(t)\steps c\), then \(c\) is final or can take a step.
\end{theorem}

The theorem quantifies over arbitrary finite prefixes. Divergent programs are
permitted; the last state of any finite prefix is not stuck.

\section{Mechanization and regression examples}
\label{sec:mechanization}

The Lean 4 development \citep{demoura2021lean4} is intrinsically scoped.
Syntax carries the exact
number of available variables, contexts and stores share that scope, and
allocation changes it from \(n\) to \(n+1\). Renaming beneath binders,
lifting runtime equality, and weakening types and semantic evidence across
allocation are explicit operations in Lean. The named presentation absorbs
these operations into capture-avoiding substitution and fresh-name
conventions. Lean implements \ruleref{rc-replace} by opening a signature in a
distinguished variable; opening substitutes the related paths for all of its
occurrences. The mechanization proves
\cref{lem:conversion-structure} through an auxiliary structural
representation of conversion, used only to obtain constructor inversion.
The transport judgment takes this inversion as its first step; Lean first
normalizes the raw conversion evidence, corresponding to
\(X\Longmapsto_{\mathsf{str}}Y\), and then follows the realization constructor. The Lean
datatype also has derived congruence constructors for functions, pairs,
singletons, selections, and intervals; normalization treats them
homomorphically.
Accordingly, the mechanized one-step theorem is
heterogeneous in the store scope and records that its result type is the
weakening of the source type; modulo renaming, it is
\cref{thm:preservation}.

The paper also suppresses the mechanization's two-valued index distinguishing
proper members from interval members. The index rules out mismatched
signatures and definitions; no proof argument depends on its value.
The deferred function coercions and pair-member closures in
\cref{fig:semantic-judgments} correspond to separate Lean datatypes.
Function codomains support identity, composition, conversion, narrowing, and
a retained source derivation. Pair members have the retained-source case
generated by \ruleref{s-pair}. Both accept a location realizing the binder
type and produce a closed coercion between the opened signatures, following
the separate judgments in \cref{fig:coercion-operations}. Lean implements
the evidence annotations and evidence-application judgments as noncomputable
recursive functions over source derivations. The paper presents their graphs
as output-annotated rules and leaves the source proof terms implicit.

The proof contains the complete static semantics and CK machine from
\cref{sec:calculus}; it imports no earlier \lambdap{} development. The main
formal result is the specialization of \cref{thm:safety} to declarative
typing. Auxiliary files establish scoped renaming and substitution,
determinism and weakening of resolution, conversion under binders,
weakening of all semantic evidence across allocation, and the store-order
lemmas used by coercion action. These lemmas support the paper-level results
without changing their statements.

Two subtyping regressions exercise the general pair rule with a
non-singleton first component:
\[
  \pairty{x}{\top}{a}{\sing{x}}
  <:
  \pairty{x}{\top}{a}{\top}
\]
and
\[
  \pairty{x}{\top}{A}{\sing{x}..\sing{x}}
  <:
  \pairty{x}{\top}{A}{\bot..\top}.
\]
The first uses \ruleref{s-pair}, reflexivity for the first component, and
\ruleref{s-top} for the member. The second uses \ruleref{s-bounds} under the
pair binder; its lower, upper, and consistency premises are respectively
\ruleref{s-bot}, \ruleref{s-top}, and reflexivity.

The interval case also has a closed end-to-end instance. Bind the identity
function at \(\top\), store the type singleton of that location in a type
pair with exact dependent interval \(\sing{x}..\sing{x}\), and subsume the
pair to the interval \(\bot..\top\). The mechanization constructs its typing
derivation and specializes \cref{thm:safety} to every finite execution from
that term.

\section{Related work}
\label{sec:related}

\paragraph{Dependent object types.}
Foundational DOT systems account for abstract type members, dependent method
types, recursive objects, and rich subtyping
\citep{amin2014foundations,rompf2016dot}. Subsequent proofs factor the
metatheory in different ways, including inert or tight typing
\citep{rapoport2017simple} and definitional-interpreter techniques
\citep{amin2017definitional}. \lambdap{} removes recursive objects,
intersections, and unions and studies paths built from immutable dependent
pairs. The resulting runtime makes allocation order available as a
well-founded measure for the evidence generated by dependent-pair subtyping.

pDOT supports types depending on paths of arbitrary length and uses singleton
types to track path equality \citep{rapoport2019pdot}. Its soundness proof
adapts syntactic DOT techniques and imposes conditions needed for path
lookup. \lambdap{} likewise separates singleton identity from abstract type
selection, while representing record traversal by pair lookup and proving
typed path resolution through store-local realization.

gDOT gives a step-indexed logical-relations model in Iris for a guarded DOT
calculus with recursive types and objects
\citep{giarrusso2020gdot,jung2018iris}. Its surface calculus exposes a later
type former and an explicit coercion term to control guarded recursive
reasoning. The finite \lambdap{} core omits recursion; coercion evidence and
the referent stratum induced by allocation order suffice for well-foundedness.

\paragraph{Coercion interpretations.}
Coercion semantics translates subtyping or inheritance derivations into
explicit evidence, making transitivity composition and raising the question
of coherence between multiple derivations
\citep{breazutannen1991coercion,crary2000inclusive}. System FC makes type
equality evidence and casts explicit in a typed intermediate language and
proves erasure before execution \citep{sulzmann2007systemfc}. Our semantic
coercions follow the derivation-directed view of evidence and preserve
realization at a fixed referent within a store. Type safety requires each interpreted
derivation to preserve realization; it does not require different derivations of the same
subtyping judgment to produce equal evidence. The dependent cases retain
source derivations in closures and interpret them only after a concrete
location is supplied.

Stratified semantic models have long used worlds or store levels to break
circular definitions involving higher-order state
\citep{ahmed2002stratified}. Location referents are ranked by their allocation
position in the immutable store, while type referents have rank zero. The
rank is used solely to justify recursive coercion action at referents
stored inside a pair.

\paragraph{Intervals, GADTs, and compilation to DOT.}
Type intervals also support a uniform account of higher-order bounded
abstraction and translucent definitions in
\(F^{\omega}_{..}\) \citep{stucki2021intervals}. \lambdap{} uses first-order
intervals as signatures of abstract pair members and retains explicit
nonemptiness premises in selection, interval formation, and interval
subtyping.

Work on GADT reasoning in Scala and cDOT studies equality and subtyping
constraints recovered by pattern matching
\citep{parreaux2019gadt,boruchgruszecki2022cdot}. Those calculi address
features absent from \lambdap{}, but they reinforce the importance of
tracking the evidence behind path-dependent type refinement. Type-preserving
translations from class-based calculi to DOT address the complementary
question of connecting a source language to a sound path-dependent core
\citep{martres2023compilation}.

\section{Conclusion}

The general dependent-pair rule is sound when its member premise remains
available until a run-time pair supplies the first-component location.
Store-indexed semantic coercions provide the required proof object. They
compose along transitivity, retain the evidence behind singleton and
selection rules, and transform realizations of referents without extending the
program syntax. Pair coercion action instantiates its member closure at the
stored component and terminates by descending in the referent rank induced by
the allocation order.

This construction proves progress and preservation for \lambdap{} with
arbitrary paths, abstract interval members, primitive transitivity, and
covariance of both pair components. It also identifies a compact proof
interface for future extensions: typed resolution, realization at a referent,
derivation-directed coercion construction, and a machine invariant whose
subsumption cases are normalized to trailing coercions.

\bibliographystyle{ACM-Reference-Format}
\bibliography{../references}

\end{document}
